%% 7-5 逆行列で連立方程式を解く（第14回：組 V7-3。ページ制限なし。偶数ページ）
%% ほかの巻・物理の本との約束：
%%   7-1：A x = b は「数の方程式 2x=6 とそっくり」、行列の「わり算」にあたるもの（逆行列）は 7-5。3つのばね（k=100 N/m、W=0.98 N）の答え 0.0294・0.0490・0.0588 m は 7-5 で。
%%   7-1 のコラム『九章算術』の方程 → 7-5 の掃き出し法と同じ考え。7-2：ななめの物差し (5,5) から s=2、t=1 を逆行列で。ヤコビ行列を逆向きに使う（手の速度→関節の角速度）。
%%   3-6：くだもの屋（120円・60円）、回路（9 V・6 V、2 Ω、2 A・0.5 A）、3つの文字（x+y+z=9…、4・3・2）を、逆行列と掃き出し法で回収。
%%   物理 力-5：A x = b は x = A^{-1} b。表計算「空いた列のたて6マスをえらび =MMULT(MINVERSE(A2:F7),G2:G7)」、np.linalg.solve。正三角形のトラスの答え -5.77、-5.77、2.89、0、5、5。
%%   物理 力-15：左の行列の行列式は1で0でないので逆行列がある。物理 p1-05 の専門書：「正方かつ正則」は「たて・横同じ大きさで、逆行列がある」こと。p1-15：階数は大学の線形代数で。
%%   3-16：最小二乗法 y=ax+b は、2つを同時に決める連立方程式（行列、7-5）。委員会の決定8：偏微分を使わず、ずれの矢印が直交する図で A^T A x = A^T b を導く（★4）。ソラさんのデータで a=1.489、b=0.032。
%%   数はすべて Python（numpy 2.5）で確かめた（v7-3/py/c5.py・out5.py）。
\usetikzlibrary{matrix}
\makeatletter
% 本文の小見出し（太字の1行）
\newcommand\kviifhead[1]{\par\addvspace{6pt}\Needspace{6\baselineskip}\noindent{\sffamily\bfseries\color{\bkTT}#1}\par\nobreak\vspace{2pt}}
% 大事な式のわく  \kviifframe{中身}
\newcommand\kviifframe[1]{\par\addvspace{4pt}{\centering\tikz\node[draw=\bkTT, line width=0.9pt, fill=bkTheme!10, rounded corners=1.5mm, inner xsep=6mm, inner ysep=2.5mm, align=center]{#1};\par}\addvspace{4pt}\noindent\ignorespaces}
% 例題の途中式の行（2.6字下げ。折り返しも同じ位置から）
\newcommand\kviifq{\par\vspace{1pt}\noindent\hangindent=2.6\zw\hangafter=0\relax}
% 小さな行列（図の説明・こたえの中だけで使う）
\newcommand\kviifsm[1]{\left(\begin{smallmatrix}#1\end{smallmatrix}\right)}
% 係数と右辺を並べた表（たて線つき）  \kviifaug{rr|r}{中身}
\newcommand\kviifaug[2]{\left(\begin{array}{#1}#2\end{array}\right)}
% 方眼と軸（kata7 の型）
\newcommand\kviifgrid[4]{%
  \draw[bkLine, line width=0.4pt, step=1] (#1,#3) grid (#2,#4);
  \draw[bkInk, line width=0.7pt, -{Stealth[length=2.2mm]}] (#1,0) -- ({#2+0.7},0) node[right, font=\footnotesize, inner sep=1pt] {$x$};
  \draw[bkInk, line width=0.7pt, -{Stealth[length=2.2mm]}] (0,#3) -- (0,{#4+0.7}) node[above, font=\footnotesize, inner sep=1pt] {$y$};
  \node[below left, font=\sffamily\footnotesize, inner sep=1.2pt] at (0,0) {O};}
% もとの家（灰色の点線）と、行列をかけた家（kata7 の型）
\newcommand\kviifhousezero{%
  \filldraw[fill=bkGray!12, draw=bkGray, line width=0.8pt, dashed, line join=round] (0,0) -- (2,0) -- (2,2) -- (1,3) -- (0,2) -- cycle;
  \draw[bkGray, line width=0.6pt, dashed] (1.25,0) -- (1.25,1) -- (1.75,1) -- (1.75,0);}
\newcommand\kviifhouse[5][bkOrange]{%
  \begin{scope}[cm={#2,#4,#3,#5,(0,0)}]
    \filldraw[fill=#1!25, draw=#1!80!black, line width=1.1pt, line join=round] (0,0) -- (2,0) -- (2,2) -- (1,3) -- (0,2) -- cycle;
    \filldraw[fill=white, draw=#1!80!black, line width=0.8pt] (1.25,0) -- (1.75,0) -- (1.75,1) -- (1.25,1) -- cycle;
  \end{scope}}
% 1列め（赤）と2列め（緑）の矢印
\newcommand\kviifbasis[4]{%
  \draw[bkRed, line width=1.3pt, -{Stealth[length=2.4mm]}] (0,0) -- (#1,#3);
  \draw[bkGreen!60!black, line width=1.3pt, -{Stealth[length=2.4mm]}] (0,0) -- (#2,#4);}
% 表計算の数式バー（kata7 の型）  \kviiffx{D2}{=MMULT(…)}
\newcommand\kviiffx[2]{\tikz[baseline=(f.base)]{\node[draw=bkLine, line width=0.6pt, fill=white, minimum height=5.5mm, inner xsep=2mm, font=\sffamily\footnotesize, minimum width=12mm] (c) {#1};
  \node[anchor=west, draw=bkLine, line width=0.6pt, fill=white, minimum height=5.5mm, inner xsep=2mm, minimum width=50mm, text width=66mm, align=left, font=\ttfamily\footnotesize] (f) at ([xshift=1mm]c.east) {#2};}}
\colorlet{kviifM}{bkTheme!12}\colorlet{kviifV}{bkGreen!14}\colorlet{kviifR}{bkYellow!35}
\makeatother

\bkchapter{7}{5}{逆行列で\\連立方程式を解く}{\bklevel{4}}{逆行列で連立方程式を解く}
\bkmichishirube{3-6,3-16,4-11,4-12,7-1,7-3,7-4}{7-9,7-11,9-10}

\begin{goal}
\begin{itemize}[label={\color{\bkTT}●}, leftmargin=1.6\zw, itemsep=0pt, topsep=1pt]
\item \bklevel{4}\enspace 2×2 の逆行列を作って確かめ、連立方程式を $\vec x=A^{-1}\vec b$ で解ける。
\item \bklevel{4}\enspace 掃き出し法で、3つの文字の連立方程式と逆行列を、表の行の操作で求められる。
\item \bklevel{4}\enspace 表計算・Python で連立方程式を解き、最小二乗法の式 $A^{\mathsf T}A\vec x=A^{\mathsf T}\vec b$ で直線を求められる。
\end{itemize}
\end{goal}

\begin{hatena}
\noindent\begin{minipage}[c]{0.5\linewidth}\raggedright
3-6では、2つの式の連立方程式を、加減法で解きました。

橋のトラスでは、棒がつながる点（節点）の1つ1つで、横とたての力のつり合いの式が2つずつできます。節点が50個なら、式は100個です。大きな電気の回路でも同じです。

式が100個もある連立方程式を、どうやって解くのでしょう。
\end{minipage}\hfill
\begin{minipage}[c]{0.46\linewidth}\centering
\begin{tikzpicture}[x=0.82mm, y=1mm]
  \fill[bkGray!35] (-3,-3) rectangle (63,-1.2);
  \draw[bkInk!80, line width=1pt] (0,0) -- (60,0);
  \draw[bkInk!80, line width=1pt] (7.5,12) -- (52.5,12);
  \foreach \kviifx in {0,15,30,45} {\draw[bkInk!80, line width=0.8pt] (\kviifx,0) -- (\kviifx+7.5,12) -- (\kviifx+15,0);}
  \foreach \kviifx in {0,15,30,45,60} {\fill[bkRed] (\kviifx,0) circle (1);}
  \foreach \kviifx in {7.5,22.5,37.5,52.5} {\fill[bkRed] (\kviifx,12) circle (1);}
  \foreach \kviifx in {15,30,45} {\draw[bk arrow, bkPurple, line width=0.9pt] (\kviifx,-1.5) -- (\kviifx,-9);}
  \node[font=\sffamily\footnotesize, anchor=south] at (30,14) {節点（赤い点）1つに式が2つ};
  \node[font=\sffamily\footnotesize, anchor=north] at (30,-10) {節点が9個 → 式は18個};
\end{tikzpicture}
\end{minipage}
\end{hatena}

\noindent 7-1では、連立方程式を $A\vec x=\vec b$ と、たった1つの式で書きました。数の方程式 $2x=6$ とそっくりの形です。$2x=6$ は、両辺を2でわって $x=3$ と解きました。行列でも、同じように「わって」解けないでしょうか。

行列には、わり算はありません。そのかわり、7-4で「つぶさない行列（行列式が0でない行列）は、元にもどせる」ことを学びました。この章では、行列の変換を\textbf{元にもどす行列}を作り、それをかけて連立方程式を解きます。そして、コンピュータが、式が100個でも100万個でも、同じ手順で解いているしくみを見ます。最後に、3-16の最小二乗法を、行列で解きます。

%% ---------------- やってみよう・思い出そう ----------------
\begin{yattemiyou}[表計算で連立方程式を解こう]
3-6の回路（\qty{9}{V} と \qty{6}{V} の電源、\qty{2}{\ohm} の抵抗3つ）の式 $4I_1+2I_2=9$、$2I_1+4I_2=6$ を、表計算の関数で解きます。
\begin{enumerate}[label={\bkcircled[bkGreen]{\arabic*}}, leftmargin=2.4\zw, itemsep=1pt, topsep=3pt]
\item A2:B3 に係数、C2:C3 に右辺（電源の電圧）を入れよう（7-1の並べ方）。
\item D2 に \texttt{=MMULT(MINVERSE(A2:B3),C2:C3)} と入れて、7-1と同じように確定しよう。D2:D3 に $I_1$、$I_2$ が出るよ。
\item 答えを元の式に入れて確かめよう。E2 に \texttt{=MMULT(A2:B3,D2:D3)} と入れると、右辺の9と6にもどるかな？
\item C2:C3 を12と0に変えてみよう。答えはどうなるかな？
\end{enumerate}
\begin{bkcenter}
\kviiffx{D2}{=MMULT(MINVERSE(A2:B3),C2:C3)}\par\vspace{1mm}
\bksheet{5}{16}{%
1 & \SetCell[c=2]{c} 係数 $A$ & & \SetCell{c} 右辺 $\vec b$ & \SetCell{c} 答え $\vec x$ & \SetCell{c} 確かめ $A\vec x$ \\
2 & \SetCell{bg=kviifM} 4 & \SetCell{bg=kviifM} 2 & \SetCell{bg=kviifV} 9 & \SetCell{bg=kviifR} 2 & 9 \\
3 & \SetCell{bg=kviifM} 2 & \SetCell{bg=kviifM} 4 & \SetCell{bg=kviifV} 6 & \SetCell{bg=kviifR} 0.5 & 6 \\}
\bktabcaption{\texttt{MINVERSE} は逆行列を作る関数（①で学ぶ）。答えは $I_1=\qty{2}{A}$、$I_2=\qty{0.5}{A}$ で、3-6と同じ。E列で元の式にもどしても、9と6になる。}
\end{bkcenter}
\end{yattemiyou}

\noindent ④で右辺を12と0にすると、答えは $I_1=4$、$I_2=-2$ になります。$I_2$ のマイナスは、右の電源の電流が、3-6の図の矢印と反対向きに流れるという意味です（3-6で決めた向きを基準にしている）。右辺を変えるだけで、すぐに答えが変わりました。表計算は、どんな計算をしているのでしょう。

\begin{omoidasou}
\bkrecall{3-6}{加減法：式に数をかけたり、式どうしをひいたりして、文字を消す。係数の表の行どうしをひいても同じ（行列の種）。}
\bkrecall{7-1}{連立方程式は $A\vec x=\vec b$（$A$ は係数行列）と書ける。2行3列の行列と3行1列の行列はかけられる（左の列の数＝右の行の数）。}
\bkrecall{7-2}{行列は「入れると出てくる機械」。1列めは $(1,\ 0)$ の行き先、2列めは $(0,\ 1)$ の行き先。}
\bkrecall{7-3}{$AB\vec x$ は $B$ を先にかける。$(AB)C=A(BC)$。単位行列 $I$ は $AI=IA=A$。}
\bkrecall{7-4}{2×2 の行列式は $ad-bc$（たすきにかけてひく）。0なら平らにつぶれて元にもどせない。$\det(AB)=\det A\times\det B$。}
\bkrecall{4-12}{内積 $\vec a\cdot\vec b=a_xb_x+a_yb_y$。直角なら内積は0。$\vec a\cdot\vec a=|\vec a|^2$。}
\bkrecall{3-16}{最小二乗法：ずれの2乗の和 $S$ をいちばん小さくする。原点を通る直線 $y=ax$ なら、$a=$（$xy$ の和）÷（$x^2$ の和）。}
\end{omoidasou}

%% ---------------- 数学の道具 ----------------
\bksection{doguu}{数学の道具}

\bksubsection{1}{元にもどす機械：逆行列}

\leavevmode\label{bk:7-5:gyaku}数の方程式 $2x=6$ を、もう一度ていねいに見ます。「2でわる」は、「$\dfrac{1}{2}$ をかける」ことです。$\dfrac{1}{2}$ は、2にかけると1になる数です。
\begin{bkhenkei}[0.4]
2x=6 & 問題の式 \\
\frac{1}{2}\times2x=\frac{1}{2}\times6 & 両辺に、左から $\frac{1}{2}$ をかける \\
1\times x=3 & $\frac{1}{2}\times2=1$（「2倍」を「元にもどす」） \\
x=3 & 1をかけても変わらない \\
\end{bkhenkei}
行列でも、同じことをしたいのです。そのためには、「行列 $A$ にかけると単位行列 $I$ になる行列」がほしくなります。単位行列 $I$ は、数の1と同じく、かけても何も変えない行列でした（7-3）。

\kviifhead{絵を元にもどす}
7-2で、家の絵に主役の行列 $A$ をかけると、ななめにのびた家になりました。この家を、元の家にもどす行列はあるでしょうか。
\begin{bkfigure}
\begin{tikzpicture}[x=4.6mm, y=4.6mm]
  \begin{scope}
    \kviifgrid{-1}{3}{-1}{4}
    \kviifhouse[bkGray]{1}{0}{0}{1}
    \node[font=\sffamily\footnotesize, anchor=north] at (1,-1.6) {元の家};
  \end{scope}
  \draw[bk arrow, bkOrange!80!black, line width=1pt] (4.6,4.2) -- (8.4,4.2);
  \node[font=\footnotesize, anchor=south] at (6.5,4.4) {$A$ をかける};
  \draw[bk arrow, \bkTT, line width=1pt] (8.4,2.0) -- (4.6,2.0);
  \node[font=\footnotesize, anchor=north] at (6.5,1.8) {$A^{-1}$ をかける};
  \begin{scope}[shift={(11,0)}]
    \kviifgrid{-1}{7}{-1}{8}
    \kviifhousezero
    \kviifhouse{2}{1}{1}{2}
    \node[font=\sffamily\footnotesize, anchor=north] at (3,-1.6) {$A$ をかけた家（面積3倍）};
  \end{scope}
\end{tikzpicture}
\bkfigcaption{主役の行列 $A$ で、家はななめにのびる（オレンジ）。元にもどす行列 $A^{-1}$ をかけると、元の家（灰色）にもどる。}\label{bk:7-5:fig-ie}
\end{bkfigure}
7-4で学んだように、主役の行列の行列式は3で、0ではありません。家はつぶれずに、面積が3倍の家になっただけです。つぶれていないので、元にもどすことができそうです。

\kviifhead{元にもどす行列を、さがしてみる}
元にもどす行列を $X$ として、$AX=I$ になる $X$ をさがします。行列×行列は「右の行列の列ごとに、行列×ベクトルをする」のでした（7-1、7-3）。$I$ の1列めは $(1,\ 0)$、2列めは $(0,\ 1)$ です。だから、$X$ の1列めは $A\vec x=\vec e_x$ の答え、2列めは $A\vec x=\vec e_y$ の答えです。2つの連立方程式を、加減法で解きます（$\vec x$ の成分を $x$、$y$ とする）。
\begin{bkhenkei}[0.5]
2x+y=1\ \cdots\text{①} & 1列め：$A\vec x=\vec e_x$ \\
x+2y=0\ \cdots\text{②} & \\
3x=2,\quad x=\tfrac{2}{3} & ①$\times2-{}$② で、$y$ が消える \\
3y=-1,\quad y=-\tfrac{1}{3} & ②$\times2-{}$① で、$x$ が消える \\
2x+y=0\ \cdots\text{①} & 2列め：$A\vec x=\vec e_y$ \\
x+2y=1\ \cdots\text{②} & \\
3x=-1,\quad x=-\tfrac{1}{3} & ①$\times2-{}$② で、$y$ が消える \\
3y=2,\quad y=\tfrac{2}{3} & ②$\times2-{}$① で、$x$ が消える \\
\end{bkhenkei}
2つの答えを列にしてならべると、元にもどす行列が見つかります。
\[ X=\begin{pmatrix}\frac{2}{3}&-\frac{1}{3}\\[2pt]-\frac{1}{3}&\frac{2}{3}\end{pmatrix}=\frac{1}{3}\begin{pmatrix}2&-1\\-1&2\end{pmatrix} \]
どちらの答えにも、3でわる形が出てきました。この3は、加減法で $x$ や $y$ の前に出てきた数 $2\times2-1\times1$、つまり行列式 $\det A$ の値です。

\noindent 見つけた行列を、こんどは $A$ に左からかけて、確かめてみます。4つの成分は、行と列の内積です（7-3）。
\[ \begin{pmatrix}2&-1\\-1&2\end{pmatrix}\begin{pmatrix}2&1\\1&2\end{pmatrix}=\begin{pmatrix}3&0\\0&3\end{pmatrix} \]
\[ \text{1行め：}\ 2\times2+(-1)\times1=3,\quad 2\times1+(-1)\times2=0 \]
\[ \text{2行め：}\ (-1)\times2+2\times1=0,\quad (-1)\times1+2\times2=3 \]
ななめ以外が0になって、$3I$（単位行列の3倍）です。$\dfrac{1}{3}$ 倍すれば、ちょうど $I$ になります。左からかけても、元にもどす行列になっていました。

\kviifhead{どんな2×2 の行列でも、同じ形になる}
文字でも、同じように解きます。$A$ の成分を $a$、$b$、$c$、$d$ とします（1行め $(a,\ b)$、2行め $(c,\ d)$）。
\begin{bkhenkei}[0.5]
ax+by=1\ \cdots\text{①} & 1列め：$A\vec x=\vec e_x$ \\
cx+dy=0\ \cdots\text{②} & \\
(ad-bc)x=d & ①$\times d-{}$②$\times b$ で、$y$ が消える \\
(ad-bc)y=-c & ②$\times a-{}$①$\times c$ で、$x$ が消える \\
ax+by=0\ \cdots\text{①} & 2列め：$A\vec x=\vec e_y$ \\
cx+dy=1\ \cdots\text{②} & \\
(ad-bc)x=-b & ①$\times d-{}$②$\times b$ で、$y$ が消える \\
(ad-bc)y=a & ②$\times a-{}$①$\times c$ で、$x$ が消える \\
\end{bkhenkei}
どれも $ad-bc$（行列式）でわれば答えです。1列めは $(d,\ -c)$、2列めは $(-b,\ a)$ を $ad-bc$ でわったもの。$a$ と $d$ が入れかわり、$b$ と $c$ にマイナスがついた形です。

\noindent この行列を $A$ に左からかけても $I$ になるか、4つの成分を1つずつ計算して確かめます（行と列の内積、7-3）。
\begin{bkhenkei}[0.5]
\text{1行1列め}\ \ d\times a+(-b)\times c=ad-bc & 行列式そのもの \\
\text{1行2列め}\ \ d\times b+(-b)\times d=0 & $db$ と $bd$ が消える \\
\text{2行1列め}\ \ (-c)\times a+a\times c=0 & $ca$ と $ac$ が消える \\
\text{2行2列め}\ \ (-c)\times b+a\times d=ad-bc & 行列式そのもの \\
\end{bkhenkei}
答えは、ななめが $ad-bc$、ほかが0の行列、つまり $(ad-bc)I$ です。だから、行列式 $ad-bc$ でわれば、ちょうど $I$ になります。
\kviifframe{$A=\begin{pmatrix}a&b\\c&d\end{pmatrix}$ を元にもどす行列は\\[4pt]$\dfrac{1}{ad-bc}\begin{pmatrix}d&-b\\-c&a\end{pmatrix}$\qquad（$ad-bc\neq0$）}
この行列を $A$ の\textbf{逆行列}（ぎゃくぎょうれつ）といい、$A^{-1}$ と書きます。\bkidx{ぎゃくぎょうれつ}{逆行列}「エー・インバース」と読みます（inverse は「逆の」という意味）。作り方は3つの手順です。
\begin{enumerate}[label={\bkcircled[\bkTF]{\arabic*}}, leftmargin=2.4\zw, itemsep=1pt, topsep=2pt]
\item 左上 $a$ と右下 $d$ を入れかえる。
\item 右上 $b$ と左下 $c$ に、マイナスをつける（場所はそのまま）。
\item 全体を、行列式 $ad-bc$ でわる（$\dfrac{1}{ad-bc}$ 倍する）。
\end{enumerate}
\noindent 主役の行列なら、次のとおりです。
\[ A^{-1}=\frac{1}{3}\begin{pmatrix}2&-1\\-1&2\end{pmatrix}=\begin{pmatrix}\dfrac{2}{3}&-\dfrac{1}{3}\\[6pt]-\dfrac{1}{3}&\dfrac{2}{3}\end{pmatrix} \]

\kviifhead{逆から見ても、図で見ても}
逆行列の1列めは、$(1,\ 0)$ の行き先です（7-2）。主役の $A^{-1}$ の1列めは $\left(\dfrac{2}{3},\ -\dfrac{1}{3}\right)$ です。これに $A$ をかけると、
\[ A\begin{pmatrix}\frac{2}{3}\\[2pt]-\frac{1}{3}\end{pmatrix}=\begin{pmatrix}\frac{4}{3}-\frac{1}{3}\\[2pt]\frac{2}{3}-\frac{2}{3}\end{pmatrix}=\begin{pmatrix}1\\0\end{pmatrix} \]
で、$(1,\ 0)$ にもどります。$A^{-1}$ は、「$A$ をかけると $(1,\ 0)$ になる点」に、$(1,\ 0)$ を送る機械なのです。
\begin{bkfigure}
\begin{tikzpicture}[x=15mm, y=15mm]
  \draw[bkLine, line width=0.4pt, step=0.3333] (-0.7,-0.7) grid (1.35,1.35);
  \draw[bkInk, line width=0.7pt, -{Stealth[length=2.2mm]}] (-0.7,0) -- (1.5,0) node[right, font=\footnotesize, inner sep=1pt] {$x$};
  \draw[bkInk, line width=0.7pt, -{Stealth[length=2.2mm]}] (0,-0.7) -- (0,1.5) node[above, font=\footnotesize, inner sep=1pt] {$y$};
  \draw[bkGray, line width=0.7pt, dashed] (0,0) rectangle (1,1);
  \filldraw[fill=bkTheme!22, draw=\bkTT, line width=1pt] (0,0) -- (0.6667,-0.3333) -- (0.3333,0.3333) -- (-0.3333,0.6667) -- cycle;
  \kviifbasis{0.6667}{-0.3333}{-0.3333}{0.6667}
  \node[below left, font=\sffamily\footnotesize, inner sep=1.2pt] at (0,0) {O};
  \node[font=\footnotesize, anchor=north west, fill=white, inner sep=1pt] at (0.68,-0.36) {$\left(\frac{2}{3},\ -\frac{1}{3}\right)$};
  \node[font=\footnotesize, anchor=south east, fill=white, inner sep=1pt] at (-0.36,0.7) {$\left(-\frac{1}{3},\ \frac{2}{3}\right)$};
  \node[font=\sffamily\footnotesize, anchor=south west, fill=white, inner sep=1pt] at (1.02,1.02) {面積1};
\end{tikzpicture}
\bkfigcaption{$A^{-1}$ の1列め（赤）と2列め（緑）。2本がつくる平行四辺形の面積は $\frac{1}{3}$（点線の正方形は面積1）。}\label{bk:7-5:fig-gyaku}
\end{bkfigure}
また、$A^{-1}$ の2本の列がつくる平行四辺形（図\ref{bk:7-5:fig-gyaku}）の面積は $\dfrac{1}{3}$ です。$A$ が面積を3倍にするので、元にもどす $A^{-1}$ は面積を $\dfrac{1}{3}$ 倍にします。7-4の $\det(AB)=\det A\times\det B$ から、$\det(A^{-1})\times\det A=\det I=1$、つまり $\det(A^{-1})=\dfrac{1}{\det A}$ です。

\kviifhead{かける順を変えても、$I$ になる}
7-3で、$AB$ と $BA$ はふつうちがう、と学びました。でも、逆行列は、右からかけても（①で $AX=I$ として作った）、左からかけても（①で確かめた）、$I$ になります。主役の行列で、右からかけたものをもう一度見ておきます。
\[ \begin{pmatrix}2&1\\1&2\end{pmatrix}\cdot\frac{1}{3}\begin{pmatrix}2&-1\\-1&2\end{pmatrix}=\frac{1}{3}\begin{pmatrix}4-1&-2+2\\2-2&-1+4\end{pmatrix}=\frac{1}{3}\begin{pmatrix}3&0\\0&3\end{pmatrix}=I \]
文字で計算しても、次のようになります（成分を計算して確かめよう）。
\[ \begin{pmatrix}a&b\\c&d\end{pmatrix}\begin{pmatrix}d&-b\\-c&a\end{pmatrix}=\begin{pmatrix}ad-bc&0\\0&ad-bc\end{pmatrix}=(ad-bc)I \]
「元にもどす」と「もどしてから、もう一度かける」は、どちらも何もしなかったのと同じなのです。
\kviifframe{$A^{-1}A=AA^{-1}=I$}

\kviifhead{行列式が0なら、逆行列はない}
$\dfrac{1}{ad-bc}$ は、$ad-bc=0$ のときは計算できません（0ではわれない）。7-4で、行列式が0の行列は平面を平らにつぶし、つぶれた絵は元にもどせないことを学びました。式で見ても、図で見ても、\textbf{行列式が0の行列には、逆行列はありません}。逆行列がある正方行列を、\textbf{正則}（せいそく）な行列といいます。\bkidx{せいそく}{正則}行列式が0でないことと、正則であることは、同じことです。

\begin{seikaku}
$A^{-1}$ を $\dfrac{1}{A}$ とは書きません。行列のわり算は決めていないからです。もし $\dfrac{B}{A}$ と書くと、$A^{-1}B$（$B$ のあとに $A^{-1}$）なのか、$BA^{-1}$（$A^{-1}$ のあとに $B$）なのか、わかりません。行列ではかける順で答えが変わる（7-3）ので、いつも「$A^{-1}$ を左からかける」「右からかける」と、場所をはっきり書きます。なお、$A^{-1}$ の $-1$ は、数の\mbox{$2^{-1}=\dfrac{1}{2}$} と同じ書き方をまねたもので、「$-1$ 乗」と読むこともあります。
\end{seikaku}

\bksubsection{2}{連立方程式を、逆行列で解く}

\leavevmode\label{bk:7-5:toku}連立方程式 $A\vec x=\vec b$ の両辺に、\textbf{左から} $A^{-1}$ をかけます。$2x=6$ の両辺に $\dfrac{1}{2}$ をかけたのと同じです。
\begin{bkhenkei}[0.4]
A\vec x=\vec b & 問題の式（7-1） \\
A^{-1}(A\vec x)=A^{-1}\vec b & 両辺に、左から $A^{-1}$ をかける \\
(A^{-1}A)\vec x=A^{-1}\vec b & かける区切りは自由（7-3：$(AB)C=A(BC)$） \\
I\vec x=A^{-1}\vec b & $A^{-1}A=I$ \\
\vec x=A^{-1}\vec b & $I$ をかけても変わらない \\
\end{bkhenkei}
\kviifframe{$A\vec x=\vec b$\qquad の答えは\qquad $\vec x=A^{-1}\vec b$}

\noindent なぜ「左から」なのでしょう。$A\vec x$ では、$A$ が $\vec x$ の左にいます。$A$ を消すには、$A$ のとなり、つまり左に $A^{-1}$ を置いて $A^{-1}A=I$ にする必要があります。右からかけた $\vec bA^{-1}$ は、2行1列の $\vec b$ に2行2列の $A^{-1}$ をかけることになり、大きさが合わないので計算できません（7-1：左の列の数＝右の行の数）。

\kviifhead{3-6のくだもの屋を、逆行列で解く}
7-1で、3-6のくだもの屋の問題を行列で書きました。ソラさんはりんご2個とみかん3個で420円、ミナさんはりんご2個とみかん1個で300円でした。りんご1個 $x$ 円、みかん1個 $y$ 円として、
\[ \begin{pmatrix}2&3\\2&1\end{pmatrix}\begin{pmatrix}x\\y\end{pmatrix}=\begin{pmatrix}420\\300\end{pmatrix} \]
逆行列を3つの手順で作り、右辺にかけます。
\begin{bkhenkei}[0.48]
\det A=2\times1-3\times2=2-6=-4 & 行列式（0でないので、逆行列がある） \\
A^{-1}=\frac{1}{-4}\begin{pmatrix}1&-3\\-2&2\end{pmatrix} & 2と1を入れかえ、3と2にマイナスをつけ、$-4$ でわる \\
A^{-1}\vec b=\frac{1}{-4}\begin{pmatrix}420-900\\-840+600\end{pmatrix} & 右辺にかける。1行めは $1\times420+(-3)\times300$ \\
A^{-1}\vec b=\frac{1}{-4}\begin{pmatrix}-480\\-240\end{pmatrix}=\begin{pmatrix}120\\60\end{pmatrix} & $-4$ でわる \\
\end{bkhenkei}
りんご120円、みかん60円で、3-6と同じ答えです。元の式に入れると、$2\times120+3\times60=420$、$2\times120+1\times60=300$ で、合っています。

\kviifhead{右辺が変わっても、同じ逆行列が使える}
3-6の回路の式
\[ \begin{pmatrix}4&2\\2&4\end{pmatrix}\begin{pmatrix}I_1\\I_2\end{pmatrix}=\begin{pmatrix}9\\6\end{pmatrix} \]
では、行列式は $16-4=12$ で、
\[ A^{-1}=\frac{1}{12}\begin{pmatrix}4&-2\\-2&4\end{pmatrix} \]
です。この逆行列は、抵抗だけで決まり、電源の電圧には関係しません。だから、電源を取りかえても、同じ $A^{-1}$ に右辺をかけるだけで答えが出ます。
\begin{bkcenter}
\begin{bktable}{colspec={Q[c,m,wd=34mm] Q[c,m,wd=48mm] Q[c,m,wd=30mm]}, rowsep=1.5pt}
電源の電圧 $\vec b$ & 計算 $\frac{1}{12}(4b_1-2b_2,\ -2b_1+4b_2)$ & 電流 $(I_1,\ I_2)$ \\
$(9,\ 6)$ & $\frac{1}{12}(36-12,\ -18+24)$ & $(2,\ 0.5)$ \\
$(12,\ 0)$ & $\frac{1}{12}(48-0,\ -24+0)$ & $(4,\ -2)$ \\
$(6,\ 6)$ & $\frac{1}{12}(24-12,\ -12+24)$ & $(1,\ 1)$ \\
$(0,\ 12)$ & $\frac{1}{12}(0-24,\ 0+48)$ & $(-2,\ 4)$ \\
\end{bktable}
\bktabcaption{同じ回路で、電源の電圧（V）だけを変えたときの電流（A）。逆行列は1回作れば、何度でも使える。}
\end{bkcenter}
1行めが、やってみようと3-6の答えです。2行めが、やってみようの④の答えです。工学では、同じしくみに、いろいろな入力（電圧・荷重）をあたえて調べることがよくあります。逆行列を1回作っておくと、入力が変わるたびに連立方程式を解き直さなくてすむのです。

\kviifhead{4-11のななめの物差しと、看板のひも}
7-2で、4-11のななめの物差し $\vec a=(2,\ 1)$、$\vec b=(1,\ 3)$ で $(5,\ 5)$ を表す問題を、行列で書きました。
\[ \begin{pmatrix}2&1\\1&3\end{pmatrix}\begin{pmatrix}s\\t\end{pmatrix}=\begin{pmatrix}5\\5\end{pmatrix} \]
行列式は $2\times3-1\times1=5$ なので、
\[ \begin{pmatrix}s\\t\end{pmatrix}=\frac{1}{5}\begin{pmatrix}3&-1\\-1&2\end{pmatrix}\begin{pmatrix}5\\5\end{pmatrix}=\frac{1}{5}\begin{pmatrix}10\\5\end{pmatrix}=\begin{pmatrix}2\\1\end{pmatrix} \]
4-11と同じ $s=2$、$t=1$ です。「出てきた $(5,\ 5)$ から、入れた $(s,\ t)$ を求める」のが、逆行列の仕事です。

\noindent 4-11の看板（左のひも30°、右のひも60°、\qty{100}{N}）の行列は、7-4で行列式が $-1$ でした。逆行列は、$-1$ でわる、つまり符号を変えるだけです。
\[ A^{-1}=\frac{1}{-1}\begin{pmatrix}\sin60°&-\cos60°\\-\sin30°&-\cos30°\end{pmatrix}=\begin{pmatrix}-\sin60°&\cos60°\\ \sin30°&\cos30°\end{pmatrix} \]
\[ \begin{pmatrix}T_1\\T_2\end{pmatrix}=A^{-1}\begin{pmatrix}0\\100\end{pmatrix}=\begin{pmatrix}100\cos60°\\100\cos30°\end{pmatrix}=\begin{pmatrix}50\\86.6\end{pmatrix} \]
4-11と同じく、$T_1=\qty{50}{N}$、$T_2=\qty{86.6}{N}$ です。7-4の⑦で、ひもが水平に近いと行列式が0に近づき、力が大きくなることを見ました。逆行列は行列式でわるので、行列式が0に近いと、答えがとても大きくなるのです。

\kviifhead{ロボットを、手の速度で動かす}
\label{bk:7-5:robot}7-2の物理・工学ではこう使うで、「ほしい手の速度から、関節の角速度を求める」計算は7-5で学ぶ、と約束しました。2つの関節のロボットの腕（$L_1=\qty{0.30}{m}$、$L_2=\qty{0.25}{m}$、$\theta_1=30°$、$\theta_2=45°$）では、7-2・7-4のヤコビ行列は
\[ J=\begin{pmatrix}-0.3915&-0.2415\\0.3245&0.0647\end{pmatrix},\qquad \det J=0.0530 \]
でした（丸める前の値）。手をまっすぐ右へ \qty{0.1}{m/s} で動かしたいとき、関節の角速度 $\dot{\vec\theta}$ は、次の式の答えです。
\[ J\dot{\vec\theta}=\begin{pmatrix}0.1\\0\end{pmatrix} \]
\[ \dot{\vec\theta}=J^{-1}\begin{pmatrix}0.1\\0\end{pmatrix}=\frac{1}{0.0530}\begin{pmatrix}0.0647&0.2415\\-0.3245&-0.3915\end{pmatrix}\begin{pmatrix}0.1\\0\end{pmatrix} \]
右辺の2行めが0なので、逆行列の1列めだけが効きます。
\[ \dot{\vec\theta}=\frac{1}{0.0530}\begin{pmatrix}0.0647\times0.1\\-0.3245\times0.1\end{pmatrix}=\frac{1}{0.0530}\begin{pmatrix}0.00647\\-0.03245\end{pmatrix}=\begin{pmatrix}0.122\\-0.612\end{pmatrix} \]
肩を \qty{0.122}{rad/s} で左回りに、ひじを \qty{0.612}{rad/s} で右回りに回せばよいのです。確かめに $J$ をかけると、次のように、ほしい速度になります。
\[ -0.3915\times0.122+0.2415\times0.612=0.100 \]
\[ 0.3245\times0.122-0.0647\times0.612=0.000 \]

\noindent ひじの角が10°で腕がのびきる近くでは、$\det J=0.0130$ まで小さくなり、同じ \qty{0.1}{m/s} のために、肩 \qty{1.47}{rad/s}、ひじ \qty{-3.47}{rad/s} もの角速度がいります。7-4の特異姿勢の近くで「関節をとても速く回さなければならない」のは、逆行列が小さな行列式でわるからです。

\Needspace{14\baselineskip}
\begin{motto}[4]
\textbf{クラメルの公式：行列式だけで答えを書く}\par
2つの文字の連立方程式
\[ \begin{cases}ax+by=p\\cx+dy=q\end{cases} \]
の答えは、逆行列をかけて成分を書き出すと、
\[ x=\frac{dp-bq}{ad-bc}=\frac{\det\begin{pmatrix}p&b\\q&d\end{pmatrix}}{\det\begin{pmatrix}a&b\\c&d\end{pmatrix}},\qquad y=\frac{aq-cp}{ad-bc}=\frac{\det\begin{pmatrix}a&p\\c&q\end{pmatrix}}{\det\begin{pmatrix}a&b\\c&d\end{pmatrix}} \]
です。$x$ は「1列めを右辺に取りかえた行列の行列式」を、もとの行列式でわったもの、$y$ は「2列めを右辺に取りかえた行列の行列式」をわったものです。くだもの屋なら、$x=\dfrac{420\times1-3\times300}{-4}=\dfrac{-480}{-4}=120$ です。これを\textbf{クラメルの公式}といい、スイスのクラメルが1750年の本で示しました。3×3 でも同じ形になりますが、行列式の計算がふえるので、大きな連立方程式には掃き出し法を使います。
\end{motto}

\bksubsection{3}{2つの機械を元にもどす順番}

\leavevmode\label{bk:7-5:junban}$B$ をかけてから $A$ をかける機械 $AB$（7-3）を、元にもどすにはどうするでしょう。朝、くつ下をはいてから、くつをはきます。夜、ぬぐときは、くつが先で、くつ下があとです。\textbf{あとからしたことを、先にもどします}。
\begin{bkfigure}
\begin{tikzpicture}[x=1mm, y=1mm]
  \node[bk node, draw=bkOrange!80!black, fill=bkOrange!12, minimum width=30mm, align=center, font=\sffamily\footnotesize] (a) at (0,12) {くつ下をはく\\（$B$）};
  \node[bk node, draw=bkOrange!80!black, fill=bkOrange!12, minimum width=30mm, align=center, font=\sffamily\footnotesize] (b) at (42,12) {くつをはく\\（$A$）};
  \node[bk node, draw=\bkTT, fill=bkTheme!12, minimum width=30mm, align=center, font=\sffamily\footnotesize] (c) at (42,-10) {くつをぬぐ\\（$A^{-1}$）};
  \node[bk node, draw=\bkTT, fill=bkTheme!12, minimum width=30mm, align=center, font=\sffamily\footnotesize] (d) at (84,-10) {くつ下をぬぐ\\（$B^{-1}$）};
  \draw[bk link] (a) -- (b);
  \draw[bk link] (b) -- (c);
  \draw[bk link] (c) -- (d);
  \node[font=\sffamily\footnotesize, anchor=west] at (62,12) {はく：$AB$（$B$ が先）};
  \node[font=\sffamily\footnotesize, anchor=east] at (22,-10) {ぬぐ：$B^{-1}A^{-1}$（$A^{-1}$ が先）};
\end{tikzpicture}
\bkfigcaption{はくのは $B$ が先（$AB$）。ぬぐのは $A^{-1}$ が先なので、右から先にかける書き方では $B^{-1}A^{-1}$。}
\end{bkfigure}
\kviifframe{$(AB)^{-1}=B^{-1}A^{-1}$\qquad（順番が逆になる）}
本当に元にもどるか、1行ずつ確かめます。
\begin{bkhenkei}[0.44]
(B^{-1}A^{-1})(AB) & もどす機械 $B^{-1}A^{-1}$ を、$AB$ のあとにかける \\
=B^{-1}(A^{-1}A)B & かける区切りは自由（7-3） \\
=B^{-1}IB=B^{-1}B & $A^{-1}A=I$、$I$ は何も変えない \\
=I & $B^{-1}B=I$ \\
\end{bkhenkei}
$A^{-1}B^{-1}$ の順では、まん中が $B^{-1}A$ になって消えません。数で確かめます。主役の行列 $A$ と、ずらしの行列 $B$ で計算します。
\[ B=\begin{pmatrix}1&1\\0&1\end{pmatrix},\qquad B^{-1}=\begin{pmatrix}1&-1\\0&1\end{pmatrix} \]
\[ AB=\begin{pmatrix}2&3\\1&3\end{pmatrix},\qquad (AB)^{-1}=\frac{1}{3}\begin{pmatrix}3&-3\\-1&2\end{pmatrix} \]
\[ B^{-1}A^{-1}=\frac{1}{3}\begin{pmatrix}3&-3\\-1&2\end{pmatrix}\ \ \text{（同じ）},\qquad A^{-1}B^{-1}=\frac{1}{3}\begin{pmatrix}2&-3\\-1&3\end{pmatrix}\ \ \text{（ちがう）} \]

\bksubsection{4}{掃き出し法：表の行の操作で解く}

\leavevmode\label{bk:7-5:hakidashi}2×2 の逆行列は、3つの手順ですぐに作れました。でも、3×3 やもっと大きな行列では、こんなに短い作り方はありません。そこで、3-6⑥で学んだ「係数の表の行どうしをひく」加減法を、きちんとした手順にします。

\kviifhead{3つの操作}
係数と右辺を、たて線で区切って1つの表にならべます。1行が1つの式です。この表に、次の3つの操作だけをくり返します。
\begin{enumerate}[label={\sffamily\bfseries\color{\bkTT}操作\arabic*}, leftmargin=4.2\zw, itemsep=1pt, topsep=2pt]
\item ある行に、0でない数をかける（またはわる）。
\item ある行から、ほかの行の何倍かをひく（またはたす）。
\item 2つの行を入れかえる。
\end{enumerate}
どの操作も、式を加減法で変形しただけなので、答えは変わりません。また、どの操作も、逆の操作（わる・たす・もう一度入れかえる）で元の表にもどせます。だから、答えがふえたり、へったりもしないのです。

\kviifhead{3つの文字の連立方程式を解く}
3-6のもっと知りたいの連立方程式を解きます。
\[ \begin{cases} x+y+z=9 & \cdots\text{①}\\ x-y=1 & \cdots\text{②}\\ y-z=1 & \cdots\text{③} \end{cases}\qquad\longrightarrow\qquad \kviifaug{rrr|r}{1&1&1&9\\1&-1&0&1\\0&1&-1&1} \]
②の $x-y=1$ は $1x+(-1)y+0z=1$、③の $y-z=1$ は $0x+1y+(-1)z=1$ と読んで、0もきちんと書きます。めざすのは、左の3列が「ななめが1、ほかが0」（単位行列 $I$）になった表です。そうなれば、各行は「$x=○$」「$y=○$」「$z=○$」と読めます。1列めから順に、ななめの数を1にし、その列のほかの数を0にしていきます。矢印の上の①②③は、表の1行め・2行め・3行めのことで、「②－①」は「2行めから1行めをひく」という意味です。
\[ \xrightarrow{\text{②}-\text{①}}\ \kviifaug{rrr|r}{1&1&1&9\\0&-2&-1&-8\\0&1&-1&1}\quad \xrightarrow{\text{②}\div(-2)}\ \kviifaug{rrr|r}{1&1&1&9\\0&1&0.5&4\\0&1&-1&1} \]
1列めは、1行め以外がもう0になりました（③の1列めは、はじめから0）。②から①をひくと「$0,\ -2,\ -1,\ -8$」です（$1-1=0$、$-1-1=-2$、$0-1=-1$、$1-9=-8$）。右辺もいっしょにひくことを忘れないようにします。次に、2列めのななめを1にしたので、2列めのほかの数を0にします。
\[ \xrightarrow[\text{③}-\text{②}]{\text{①}-\text{②}}\ \kviifaug{rrr|r}{1&0&0.5&5\\0&1&0.5&4\\0&0&-1.5&-3}\quad \xrightarrow{\text{③}\div(-1.5)}\ \kviifaug{rrr|r}{1&0&0.5&5\\0&1&0.5&4\\0&0&1&2} \]
最後に、3列めのほかの数を0にします。①から③の0.5倍をひき、②から③の0.5倍をひきます。
\[ \xrightarrow[\text{②}-0.5\times\text{③}]{\text{①}-0.5\times\text{③}}\ \kviifaug{rrr|r}{1&0&0&4\\0&1&0&3\\0&0&1&2} \]
左の3列が単位行列になりました。1行めは $1x+0y+0z=4$、つまり $x=4$ です。同じように $y=3$、$z=2$ で、3-6と同じ答えです。元の式に入れると、$4+3+2=9$、$4-3=1$、$3-2=1$ で合っています。

\noindent このように、1列ずつ、ななめ以外の数を0にして「はき出して」いく解き方を、\textbf{掃き出し法}（はきだしほう）といいます。\bkidx{はきだしほう}{掃き出し法}ほうきでごみを掃き出すように、じゃまな数を0にしていくので、この名前があります。7-1のコラムで見た『九章算術』の「方程」も、係数を算木で表に並べ、列（いまの行）どうしをひいて解いていました。2000年前の計算と同じ考え方です。

\kviifhead{0の所では、行を入れかえる}
ななめにしたい場所が0のときは、その数でわって1にすることができません。そのときは、その列が0でない行と入れかえます（操作3）。式を書く順番を変えるだけなので、答えは変わりません。$y=2$、$x+y=5$ で見てみます。矢印の「①$\leftrightarrow$②」は、1行めと2行めを入れかえるという意味です。
\[ \kviifaug{rr|r}{0&1&2\\1&1&5}\ \xrightarrow{\text{①}\leftrightarrow\text{②}}\ \kviifaug{rr|r}{1&1&5\\0&1&2}\ \xrightarrow{\text{①}-\text{②}}\ \kviifaug{rr|r}{1&0&3\\0&1&2} \]
はじめの表は、1行め1列めが0です（$y=2$ は $0x+1y=2$）。2行めと入れかえると、ななめは1になりました。あとは今までどおり、①－② で $x=3$、$y=2$ です。3つの文字のときも同じです（例題5）。

\begin{seikaku}
\textbf{コンピュータは、いちばん大きな数を上にする}\par
0でなくても、とても小さな数でわると、誤差が大きくなります。そこでコンピュータは、ななめにしたい場所が0でなくても、その列でいちばん大きな数の行を上に入れかえてから計算します（部分ピボット選択という）。また、コンピュータは、ななめより下だけを0にしてから（前進消去）、下の行から順に答えを代入していきます（後退代入）。こうすると計算の回数が少なくてすみます。この方法を\textbf{ガウスの消去法}といいます。
\end{seikaku}

\kviifhead{表計算で掃き出す}
掃き出し法の1回1回の操作は、表計算の「式のコピー」と相性がよいのです。たとえば「②－①」は、2行めの各マスから、1行めの同じ列のマスをひく式を、横にコピーするだけです。
\begin{bkcenter}
\kviiffx{A6}{=A3-A2}\par\vspace{1mm}
\bksheet{4}{14}{%
1 & $x$ & $y$ & $z$ & 右辺 \\
2 & 1 & 1 & 1 & 9 \\
3 & 1 & $-1$ & 0 & 1 \\
4 & 0 & 1 & $-1$ & 1 \\
5 & 1 & 1 & 1 & 9 \\
6 & \SetCell{bg=kviifR} 0 & \SetCell{bg=kviifR} $-2$ & \SetCell{bg=kviifR} $-1$ & \SetCell{bg=kviifR} $-8$ \\
7 & 0 & 1 & $-1$ & 1 \\}
\bktabcaption{2〜4行めがはじめの表、5〜7行めが「②－①」をしたあとの表。A5 には \texttt{=A2}、A6 には \texttt{=A3-A2}、A7 には \texttt{=A4} と入れて、D列まで右へコピーする。次の操作は、この下に同じように続ける。}
\end{bkcenter}

\kviifhead{掃き出し法で、逆行列を作る}
表の右に単位行列 $I$ をならべて、左が $I$ になるまで掃き出すと、右に逆行列 $A^{-1}$ が出てきます。主役の行列でやってみます。
\[ \kviifaug{rr|rr}{2&1&1&0\\1&2&0&1}\ \xrightarrow{\text{①}\div2}\ \kviifaug{rr|rr}{1&0.5&0.5&0\\1&2&0&1} \]
\[ \xrightarrow{\text{②}-\text{①}}\ \kviifaug{rr|rr}{1&0.5&0.5&0\\0&1.5&-0.5&1}\ \xrightarrow{\text{②}\div1.5}\ \kviifaug{rr|rr}{1&0.5&0.5&0\\0&1&-\frac{1}{3}&\frac{2}{3}} \]
\[ \xrightarrow{\text{①}-0.5\times\text{②}}\ \kviifaug{rr|rr}{1&0&\frac{2}{3}&-\frac{1}{3}\\0&1&-\frac{1}{3}&\frac{2}{3}} \]
右に出てきたのは、①で作った主役の逆行列 $A^{-1}$ と同じ行列です（①$-0.5\times$② の右上は $0-0.5\times\frac{2}{3}=-\frac{1}{3}$）。

\noindent なぜ逆行列が出てくるのでしょう。行の操作は、どれも「ある行列を左からかける」ことと同じです（たとえば②－①は、1行め $(1,\ 0)$、2行め $(-1,\ 1)$ の行列を左からかけること）。全部の操作をまとめた行列を $G$ とすると、左の $A$ は $GA$ に、右の $I$ は $GI=G$ になります。左が $I$ になったのなら $GA=I$ なので、$G$ は $A$ の逆行列です。そして右には、その $G$ がそのまま出ているのです。

\bksubsection{5}{コンピュータは、どう解くか}

\leavevmode\label{bk:7-5:computer}\bkhatenaref の答えです。式が100個あっても、コンピュータがすることは、掃き出し法（ガウスの消去法）の3つの操作のくり返しです。手順は、式が2個のときとまったく同じで、くり返す回数がふえるだけです。

\kviifhead{かけ算の回数}
ガウスの消去法で、$n$ 個の文字の連立方程式を解くのに必要なかけ算（わり算もふくむ）の回数は、およそ $\dfrac{n^3}{3}$ 回です。1つの列を0にするのに、およそ $n\times n$ 回の計算をし、それを $n$ 列ぶんくり返すからです（下の行ほど、残りの列が少ないので3でわる）。
\begin{bkcenter}
\begin{bktable}{colspec={Q[l,m,wd=34mm] *{4}{Q[c,m,wd=20mm]}}, rowsep=1.5pt}
式の数 $n$ & 3 & 10 & 100 & 1000 \\
かけ算の回数（約） & 9 & 333 & 33万 & 3.3億 \\
\end{bktable}
\bktabcaption{式が10倍になると、計算の回数は約1000倍になる。それでも、いまのパソコンは1秒間に何億回も計算できるので、1000個の式でも1秒もかからずに解ける。}
\end{bkcenter}

\kviifhead{表計算：MMULT(MINVERSE())}
物理第1巻（力-5）では、正三角形のトラスの部材の力と支点の反力（わからない数6つ）を、6つの式で求めました。表計算では、行列を A2:F7、右辺を G2:G7 に入れ、\textbf{空いた列のたて6マスをえらび} \texttt{=MMULT(MINVERSE(A2:F7),G2:G7)} と入れます（物理の本と同じ書き方）。
\begin{bkcenter}
\bksheet{7}{12.5}{%
1 & $N_{\mathrm{AC}}$ & $N_{\mathrm{BC}}$ & $N_{\mathrm{AB}}$ & $R_{\mathrm{A}x}$ & $R_{\mathrm{A}y}$ & $R_{\mathrm B}$ & 右辺 \\
2 & 0.5 & 0 & 1 & 1 & 0 & 0 & 0 \\
3 & 0.866 & 0 & 0 & 0 & 1 & 0 & 0 \\
4 & 0 & $-0.5$ & $-1$ & 0 & 0 & 0 & 0 \\
5 & 0 & 0.866 & 0 & 0 & 0 & 1 & 0 \\
6 & $-0.5$ & 0.5 & 0 & 0 & 0 & 0 & 0 \\
7 & $-0.866$ & $-0.866$ & 0 & 0 & 0 & 0 & 10 \\
8 & \SetCell{bg=kviifR} $-5.77$ & \SetCell{bg=kviifR} $-5.77$ & \SetCell{bg=kviifR} 2.89 & \SetCell{bg=kviifR} 0 & \SetCell{bg=kviifR} 5 & \SetCell{bg=kviifR} 5 & ← 答え \\}
\bktabcaption{物理 力-5の正三角形のトラス（荷重 \qty{10}{kN}）。8行めは、答えを横に並べ直したもの。マイナスは押される力（圧縮）。}
\end{bkcenter}
\texttt{MINVERSE} は逆行列を作る関数、\texttt{MMULT} は行列のかけ算（7-1）です。2つを組み合わせて、$A^{-1}\vec b$ を一度に計算しています。確定のしかたは、7-1と同じです。

\kviifhead{Python：np.linalg.solve}
Python では、\texttt{np.linalg.inv} で逆行列が、\texttt{np.linalg.solve} で連立方程式の答えが求められます。3-6の回路で確かめます。
\Needspace{9\baselineskip}
\begin{bkcode}[Python]
import numpy as np
A = np.array([[4, 2],
              [2, 4]])           \# 係数行列
b = np.array([9, 6])             \# 右辺
print(np.linalg.inv(A))          \# 逆行列
print(np.linalg.inv(A) @ b)      \# 逆行列×右辺
print(np.linalg.solve(A, b))     \# 連立方程式を解く（7-5）
\end{bkcode}
\begin{bkcodeout}[実行結果]
[[ 0.33333333 -0.16666667]
 [-0.16666667  0.33333333]]
[2.  0.5]
[2.  0.5]
\end{bkcodeout}
逆行列の1行めは、$\dfrac{4}{12}=0.333\cdots$、$-\dfrac{2}{12}=-0.166\cdots$ です。2つの方法で、同じ答え $(2,\ 0.5)$ が出ました。コンピュータに大きな連立方程式を解かせるときは、\texttt{solve} を使います。\texttt{solve} は逆行列を全部は作らずに、ガウスの消去法で直接答えを出すので、計算が速く、丸めのずれも小さいからです。

\kviifhead{3つのばねを、Python で解く}
7-1で、3つのばね（ばね定数 $k=\qty{100}{N/m}$、球1個の重さ \qty{0.98}{N}）のつり合いを、$K\vec u=\vec F$ と行列で書きました。答えは7-5で求める約束でした。
\[ \begin{pmatrix}2k&-k&0\\-k&2k&-k\\0&-k&k\end{pmatrix}\begin{pmatrix}u_1\\u_2\\u_3\end{pmatrix}=\begin{pmatrix}0.98\\0.98\\0.98\end{pmatrix} \]
\Needspace{11\baselineskip}
\begin{bkcode}[Python]
import numpy as np
k = 100
K = np.array([[2*k, -k, 0],
              [-k, 2*k, -k],
              [0, -k, k]])           \# ばねの行列（剛性行列）
F = np.array([0.98, 0.98, 0.98])     \# 球の重さ（N）
u = np.linalg.solve(K, F)            \# 連立方程式を解く（7-5）
print(u)                             \# のび u1, u2, u3（m）
print(K @ u)                         \# 元の式に入れて確かめる
\end{bkcode}
\begin{bkcodeout}[実行結果]
[0.0294 0.049  0.0588]
[0.98 0.98 0.98]
\end{bkcodeout}
$u_1=\qty{0.0294}{m}$、$u_2=\qty{0.0490}{m}$、$u_3=\qty{0.0588}{m}$ です（物理 力-16と同じ）。\texttt{0.049} のあとの空白は、numpy が数の位置をそろえるために入れたものです。2行めで、元の式に入れると右辺の0.98にもどることも確かめました。

\bksubsection{6}{解けないとき、解きにくいとき}

\leavevmode\label{bk:7-5:tokenai}行列式が0の連立方程式を、掃き出し法で解こうとすると、どうなるでしょう。7-4④の例で見ます。
\[ \text{（ア）}\ \begin{cases}x+2y=3\\2x+4y=6\end{cases}\qquad\qquad \text{（イ）}\ \begin{cases}x+2y=3\\2x+4y=5\end{cases} \]
どちらも、②から①の2倍をひきます。
\[ \text{（ア）}\ \kviifaug{rr|r}{1&2&3\\2&4&6}\ \xrightarrow{\text{②}-2\times\text{①}}\ \kviifaug{rr|r}{1&2&3\\0&0&0} \]
\[ \text{（イ）}\ \kviifaug{rr|r}{1&2&3\\2&4&5}\ \xrightarrow{\text{②}-2\times\text{①}}\ \kviifaug{rr|r}{1&2&3\\0&0&-1} \]
2行めの左が、どちらも全部0になりました。2列めのななめを1にできないので、掃き出しはここで止まります。
\begin{itemize}[leftmargin=1.6\zw, itemsep=2pt, topsep=2pt, label={\color{\bkTT}●}]
\item （ア）の2行めは「$0x+0y=0$」で、どんな $x$、$y$ でも成り立ちます。のこる式は $x+2y=3$ の1つだけなので、$y$ を決めると $x=3-2y$ が決まります。\textbf{答えは無数}です。
\item （イ）の2行めは「$0x+0y=-1$」、つまり $0=-1$ です。どんな $x$、$y$ でも成り立ちません。\textbf{答えはありません}。
\end{itemize}
7-4④と同じ結論です。表計算の \texttt{MINVERSE} では \texttt{\#NUM!}（計算できない）と出て、Python の \texttt{np.linalg.solve} では、次のように「Singular matrix（正則でない行列）」と出て止まります。
\begin{bkcodeout}[実行結果]
numpy.linalg.LinAlgError: Singular matrix
\end{bkcodeout}

\begin{seikaku}
（ア）では、式は2つありましたが、2つめは1つめの2倍で、新しいことを何も言っていませんでした。掃き出したあとに残る「左が全部0ではない行」の数を、行列の\textbf{階数}（かいすう）といいます。（ア）・（イ）の行列の階数は1です。階数は「本当に役に立っている式の数」です。わからない数の数より階数が小さいと、答えは1つに決まりません。物理の次元解析（力-15）で、無次元数がいくつできるかも、この階数で決まります。くわしくは大学の線形代数で学びます。
\end{seikaku}

\kviifhead{行列式が0に近いと、答えがゆれる}
行列式が0でなくても、0にとても近いと、こまったことが起きます。
\[ \begin{cases}x+y=2\\x+1.001y=2.001\end{cases}\qquad \det=1\times1.001-1\times1=0.001 \]
答えは $x=1$、$y=1$ です。ところが、右辺の2.001を、たった0.001だけ変えて2.002にすると、答えは $x=0$、$y=2$ に大きく変わります（元の式に入れて確かめよう）。2本の直線がほとんど平行なので、少し動かしただけで、交わる点が遠くへ動いてしまうのです。
\begin{bkfigure}
\begin{tikzpicture}[x=11mm, y=11mm]
  \draw[bkLine, line width=0.4pt] (-0.5,-0.5) grid (3.5,3);
  \draw[bkInk, line width=0.7pt, -{Stealth[length=2.2mm]}] (-0.5,0) -- (3.7,0) node[right, font=\footnotesize, inner sep=1pt] {$x$};
  \draw[bkInk, line width=0.7pt, -{Stealth[length=2.2mm]}] (0,-0.5) -- (0,3.2) node[above, font=\footnotesize, inner sep=1pt] {$y$};
  \node[below left, font=\sffamily\footnotesize, inner sep=1.2pt] at (0,0) {O};
  \draw[bkBlue, line width=1.3pt] (-0.4,2.4) -- (2.4,-0.4);
  \draw[bkOrange!90!black, line width=1.3pt, dashed] (-0.5,1.75) -- (3.5,-0.25);
  \draw[bkRed, line width=1.3pt, dashed] (-0.5,2.25) -- (3.5,0.25);
  \fill[bkInk] (1,1) circle (0.06) (0,2) circle (0.06);
  \node[font=\footnotesize, anchor=north east, fill=white, inner sep=1pt] at (0.95,0.92) {$(1,\ 1)$};
  \node[font=\footnotesize, anchor=south east, fill=white, inner sep=1pt] at (-0.08,2.05) {$(0,\ 2)$};
\end{tikzpicture}
\bkfigcaption{青が $x+y=2$。オレンジの点線と赤の点線は、2つめの式の右辺を少し変えた直線。ほとんど平行なので、交わる点は $(1,\ 1)$ から $(0,\ 2)$ へ大きく動く（見やすいように、かたむきの差を大きくしてかいた）。}
\end{bkfigure}
はかった値には、いつも少しのずれがあります（2-6）。行列式が0に近い連立方程式では、そのずれが答えの大きなずれになってしまいます。このような連立方程式を「条件が悪い」といいます。7-4のひもの角やロボットの特異姿勢も、同じ理由で気をつけなければならないのです。

\bksubsection{7}{最小二乗法を行列で解く}

\leavevmode\label{bk:7-5:saisho}3-16で、ソラさんのばねのデータに、ずれの2乗の和がいちばん小さくなる直線を引きました。原点を通る直線 $y=ax$ なら、$a=$（$xy$ の和）÷（$x^2$ の和）でした。3-16では、「$a$ と $b$ の2つを同時に決める $y=ax+b$ のときも、たくさんの式を同時に解けば求められます（行列、7-5）」と約束しました。ここでは、その方法を、行列とベクトルの図で作ります。
\begin{bkcenter}
\begin{bktable}{colspec={Q[l,m,wd=34mm] *{7}{Q[c,m,wd=10mm]}}, rowsep=1.5pt}
おもりの数 $x$（個） & 0 & 1 & 2 & 3 & 4 & 5 & 6 \\
のび $y$（cm） & 0 & 1.6 & 2.9 & 4.6 & 5.9 & 7.6 & 8.9 \\
\end{bktable}
\bktabcaption{3-0・3-16のソラさんのデータ}
\end{bkcenter}

\kviifhead{式が多すぎる連立方程式}
もし7つの点が全部、直線 $y=ax+b$ の上にあるなら、次の7つの式が全部成り立つはずです。
\[ \begin{pmatrix}0&1\\1&1\\2&1\\3&1\\4&1\\5&1\\6&1\end{pmatrix}\begin{pmatrix}a\\b\end{pmatrix}=\begin{pmatrix}0\\1.6\\2.9\\4.6\\5.9\\7.6\\8.9\end{pmatrix} \]
1行めは $0a+1b=0$、2行めは $1a+1b=1.6$、…です。
左の7行2列の行列を $A$、右辺を $\vec b$、わからない $(a,\ b)$ を $\vec x$ と書くと、$A\vec x=\vec b$ です。でも、わからない数2つに、式が7つもあります。点は少しずつずれているので、7つ全部を成り立たせる $a$、$b$ はありません。そこで、ずれ $\vec b-A\vec x$（7つの成分）の2乗の和 $S$ を、いちばん小さくします。$S$ は、7つの成分をもつベクトル（4-11：成分が7つでもベクトル）$\vec b-A\vec x$ の長さの2乗です（三平方の定理、4-4・4-12）。
\[ S=|\vec b-A\vec x|^2 \]

\kviifhead{いちばん小さい例：点2つ、文字1つ}
考え方を、図にかける小さな例で見ます。点 $(1,\ 2)$、$(2,\ 3)$ に、原点を通る直線 $y=ax$ を引きます。式は $1a=2$、$2a=3$ の2つで、どちらも成り立つ $a$ はありません。

\noindent 2つの $x$ を並べた矢印 $\vec u=(1,\ 2)$ と、2つの $y$ を並べた矢印 $\vec b=(2,\ 3)$ を、1つの平面にかきます（横の軸は1つめの点、たての軸は2つめの点の値）。直線の式が予想する $y$ の組は $(1a,\ 2a)=a\vec u$ で、$a$ を変えると、$\vec u$ の向きの直線の上を動きます。ずれ $\vec b-a\vec u$ の長さがいちばん短くなるのは、図\ref{bk:7-5:fig-chokkaku}のように、\textbf{ずれの矢印が、直線に直角になるとき}です（点から直線へは、垂線がいちばん短い）。
\begin{bkfigure}
\begin{tikzpicture}[x=13mm, y=13mm]
  \draw[bkLine, line width=0.4pt] (-0.3,-0.3) grid (2.6,3.6);
  \draw[bkInk, line width=0.7pt, -{Stealth[length=2.2mm]}] (-0.3,0) -- (2.8,0) node[right, font=\footnotesize, align=left, inner sep=1pt] {1つめの点の値};
  \draw[bkInk, line width=0.7pt, -{Stealth[length=2.2mm]}] (0,-0.3) -- (0,3.8) node[above, font=\footnotesize, inner sep=1pt] {2つめの点の値};
  \draw[bkGray, line width=0.6pt, dashed] (-0.15,-0.3) -- (1.8,3.6);
  \draw[\bkTT, line width=1.4pt, -{Stealth[length=2.4mm]}] (0,0) -- (1,2);
  \draw[bkRed, line width=1.4pt, -{Stealth[length=2.4mm]}] (0,0) -- (2,3);
  \draw[bkGreen!60!black, line width=1.4pt, -{Stealth[length=2.4mm]}] (1.6,3.2) -- (2,3);
  \fill[bkInk] (1.6,3.2) circle (0.05);
  \draw[bkInk, line width=0.5pt] ($(1.6,3.2)+(-0.0894,-0.1789)$) -- ++(0.1789,-0.0894) -- ++(0.0894,0.1789);
  \node[font=\footnotesize, text=\bkTT, anchor=west, fill=white, inner sep=1pt] at (1.05,1.75) {$\vec u=(1,\ 2)$};
  \node[font=\footnotesize, text=bkRed!80!black, anchor=north west, fill=white, inner sep=1pt] at (1.25,1.5) {$\vec b=(2,\ 3)$};
  \node[font=\footnotesize, anchor=east, fill=white, inner sep=1pt] at (1.5,3.35) {$1.6\vec u$};
  \node[font=\footnotesize, text=bkGreen!50!black, anchor=west, fill=white, inner sep=1pt] at (2.05,3.05) {ずれ};
\end{tikzpicture}
\bkfigcaption{$a\vec u$ は点線の直線の上を動く。$\vec b$（赤）からの距離（緑のずれ）がいちばん短いのは、ずれが直線に直角になる点 $1.6\vec u$ のとき。}\label{bk:7-5:fig-chokkaku}
\end{bkfigure}
直角なら、内積は0です（4-12）。$\vec u\cdot(\vec b-a\vec u)=0$ を解きます。
\begin{bkhenkei}[0.45]
\vec u\cdot(\vec b-a\vec u)=0 & ずれが $\vec u$ に直角 \\
\vec u\cdot\vec b-a\,\vec u\cdot\vec u=0 & 内積をかっこの中に分ける \\
a=\frac{\vec u\cdot\vec b}{\vec u\cdot\vec u} & $a$ について解く \\
a=\frac{1\times2+2\times3}{1\times1+2\times2}=\frac{8}{5}=1.6 & 成分を入れる \\
\end{bkhenkei}
$\vec u\cdot\vec b$ は「$xy$ の和」、$\vec u\cdot\vec u$ は「$x^2$ の和」です。3-16で平方完成から作った「$a$ は、$xy$ の和を $x^2$ の和でわったもの」と、まったく同じ式が、直角の図から出てきました。$a$ を少し変えて、$S=(2-a)^2+(3-2a)^2$ を表にしても確かめられます。
\begin{bkcenter}
\begin{bktable}{colspec={Q[l,m,wd=30mm] *{5}{Q[c,m,wd=13mm]}}, rowsep=1.5pt}
$a$ & 1.4 & 1.5 & 1.6 & 1.7 & 1.8 \\
$S$ & 0.40 & 0.25 & \SetCell{bg=kviifR} 0.20 & 0.25 & 0.40 \\
\end{bktable}
\bktabcaption{$a=1.6$ で $S$ がいちばん小さい（直角の図の答えと同じ）。}
\end{bkcenter}

\kviifhead{点3つ、文字2つ：ずれが「面」に直角}
次に、点 $(0,\ 1)$、$(1,\ 3)$、$(2,\ 2)$ に、直線 $y=ax+b$ を引きます。予想する $y$ の組は、
\[ \begin{pmatrix}0a+b\\1a+b\\2a+b\end{pmatrix}=a\begin{pmatrix}0\\1\\2\end{pmatrix}+b\begin{pmatrix}1\\1\\1\end{pmatrix}=A\begin{pmatrix}a\\b\end{pmatrix},\qquad A=\begin{pmatrix}0&1\\1&1\\2&1\end{pmatrix} \]
です。3つの値の組なので、3次元の矢印です。$a$ と $b$ を変えると、この矢印は、$A$ の2本の列 $(0,\ 1,\ 2)$ と $(1,\ 1,\ 1)$ がつくる\textbf{平面}の上を動きます（4-11のななめの物差しの3次元版）。はかった値 $\vec b=(1,\ 3,\ 2)$ は、この平面からはみ出しています。いちばん近い点は、ずれの矢印がこの平面に直角になる点です。
\begin{bkfigure}
\begin{tikzpicture}[x=1mm, y=1mm]
  \filldraw[fill=bkTheme!14, draw=\bkTT!70, line width=0.5pt] (0,0) -- (80,0) -- (105,28) -- (25,28) -- cycle;
  \coordinate (O) at (20,6);
  \coordinate (P) at (68,17);
  \coordinate (B) at (68,52);
  \draw[\bkTT, line width=1.2pt, -{Stealth[length=2.2mm]}] (O) -- (52,7);
  \draw[\bkTT, line width=1.2pt, -{Stealth[length=2.2mm]}] (O) -- (40,22);
  \draw[bkInk!70, line width=0.8pt, dashed] (O) -- (P);
  \draw[bkRed, line width=1.4pt, -{Stealth[length=2.4mm]}] (O) -- (B);
  \draw[bkGreen!60!black, line width=1.4pt, -{Stealth[length=2.4mm]}] (P) -- (B);
  \fill[bkInk] (O) circle (0.8) (P) circle (0.8);
  \draw[bkInk, line width=0.5pt] ($(P)!4mm!(O)$) -- ++(0,4) -- ($(P)+(0,4)$);
  \node[font=\sffamily\footnotesize, anchor=north east, inner sep=1pt] at ($(O)+(-0.5,-0.5)$) {O};
  \node[font=\footnotesize, text=\bkTT, anchor=north, inner sep=1.5pt] at (46,6.2) {$(1,\ 1,\ 1)$};
  \node[font=\footnotesize, text=\bkTT, anchor=south east, inner sep=1.5pt] at (39,22) {$(0,\ 1,\ 2)$};
  \node[font=\footnotesize, text=bkRed!80!black, anchor=east, inner sep=1.5pt] at (40,32) {$\vec b=(1,\ 3,\ 2)$};
  \node[font=\footnotesize, anchor=north west, inner sep=1.5pt] at (69,16) {$A\vec x=(1.5,\ 2,\ 2.5)$};
  \node[font=\footnotesize, text=bkGreen!50!black, anchor=west, inner sep=1.5pt] at (69.5,36) {ずれ $(-0.5,\ 1,\ -0.5)$};
  \node[font=\sffamily\footnotesize, anchor=south east, inner sep=1.5pt] at (103,1) {2本の列がつくる平面};
\end{tikzpicture}
\bkfigcaption{$A$ の2本の列（青）がつくる平面（色をぬった面）。$A\vec x$ はこの面の上を動く。はかった値 $\vec b$（赤）にいちばん近いのは、ずれ（緑）が面に直角になる点（ようすをかいた図で、長さや角は正確ではない）。}\label{bk:7-5:fig-men}
\end{bkfigure}
ずれが平面に直角ということは、平面をつくる2本の列の\textbf{どちらとも}直角ということです。内積で書くと、2つの式になります。
\[ \begin{pmatrix}0\\1\\2\end{pmatrix}\cdot(\vec b-A\vec x)=0,\qquad \begin{pmatrix}1\\1\\1\end{pmatrix}\cdot(\vec b-A\vec x)=0 \]
この2つの内積を、行列のかけ算でまとめます。内積 $\vec u\cdot\vec v$ は、$\vec u$ を転置した行ベクトルと $\vec v$ の積でした（7-1）。2本の列を転置して行にならべた行列は、$A$ の転置 $A^{\mathsf T}$ です。
\begin{bkhenkei}[0.5]
A^{\mathsf T}(\vec b-A\vec x)=\vec0 & 2つの内積が0、をまとめた（$A^{\mathsf T}$ の行は $A$ の列） \\
A^{\mathsf T}\vec b-A^{\mathsf T}A\vec x=\vec0 & かっこを開く（行列のかけ算も分けられる） \\
A^{\mathsf T}A\vec x=A^{\mathsf T}\vec b & $A^{\mathsf T}A\vec x$ を右辺へ移して、左右を入れかえた \\
\end{bkhenkei}
\kviifframe{いちばんよく合う $\vec x$ は\qquad $A^{\mathsf T}A\vec x=A^{\mathsf T}\vec b$\qquad の答え}
この式を\textbf{正規方程式}（せいきほうていしき）といいます。\bkidx{せいきほうていしき}{正規方程式}$A$ は3行2列、$A^{\mathsf T}$ は2行3列なので、$A^{\mathsf T}A$ は2行2列の行列です（7-1：かけられる大きさ）。式が多すぎた連立方程式が、ちょうど文字の数と同じ、2つの式になりました。あとは①の逆行列で解けます。
\[ A^{\mathsf T}A=\begin{pmatrix}0&1&2\\1&1&1\end{pmatrix}\begin{pmatrix}0&1\\1&1\\2&1\end{pmatrix}=\begin{pmatrix}5&3\\3&3\end{pmatrix} \]
\[ A^{\mathsf T}\vec b=\begin{pmatrix}0&1&2\\1&1&1\end{pmatrix}\begin{pmatrix}1\\3\\2\end{pmatrix}=\begin{pmatrix}7\\6\end{pmatrix} \]
成分は、行と列の内積です。たとえば、$A^{\mathsf T}A$ の1行1列めは、$0\times0+1\times1+2\times2$ で5です。正規方程式の行列式は、$5\times3-3\times3$ で6です。だから、
\[ \begin{pmatrix}a\\b\end{pmatrix}=\frac{1}{6}\begin{pmatrix}3&-3\\-3&5\end{pmatrix}\begin{pmatrix}7\\6\end{pmatrix}=\frac{1}{6}\begin{pmatrix}3\\9\end{pmatrix}=\begin{pmatrix}0.5\\1.5\end{pmatrix} \]
直線は $y=0.5x+1.5$ です。ずれは、次のとおりです。
\[ \vec b-A\vec x=(1-1.5,\ \ 3-2,\ \ 2-2.5)=(-0.5,\ \ 1,\ \ -0.5) \]
2本の列との内積は、どちらも0です。
\[ 0\times(-0.5)+1\times1+2\times(-0.5)=0,\qquad 1\times(-0.5)+1\times1+1\times(-0.5)=0 \]たしかに、どちらとも直角です。このときの $S$ は $0.25+1+0.25=1.5$ です。$a$ や $b$ を少し変えると、$S$ はどれも1.5より大きくなります。
\begin{bkcenter}
\begin{bktable}{colspec={Q[l,m,wd=26mm] *{3}{Q[c,m,wd=17mm]}}, rowsep=1.5pt}
& $b=1.4$ & $b=1.5$ & $b=1.6$ \\
$a=0.4$ & 1.64 & 1.55 & 1.52 \\
$a=0.5$ & 1.53 & \SetCell{bg=kviifR} 1.50 & 1.53 \\
$a=0.6$ & 1.52 & 1.55 & 1.64 \\
\end{bktable}
\bktabcaption{$a$ と $b$ を0.1ずつ変えたときの $S$。まん中の $a=0.5$、$b=1.5$ がいちばん小さい。}
\end{bkcenter}

\kviifhead{ソラさんのデータで}
点が7つでも、考え方は同じです。$A^{\mathsf T}A$ と $A^{\mathsf T}\vec b$ の成分は、表の和で計算できます。$A$ の1列めは $x$、2列めは1が7個なので、
\[ A^{\mathsf T}A=\begin{pmatrix}x^2\text{ の和}&x\text{ の和}\\x\text{ の和}&\text{点の数}\end{pmatrix},\qquad A^{\mathsf T}\vec b=\begin{pmatrix}xy\text{ の和}\\y\text{ の和}\end{pmatrix} \]
です。ソラさんのデータでは、$x$ の和21、$x^2$ の和91、$y$ の和31.5、$xy$ の和136.2（3-16の表）、点の数7です。
\begin{bkhenkei}[0.6]
\begin{pmatrix}91&21\\21&7\end{pmatrix}\begin{pmatrix}a\\b\end{pmatrix}=\begin{pmatrix}136.2\\31.5\end{pmatrix} & 正規方程式 \\
\det=91\times7-21\times21=196 & 行列式（$637-441$） \\
a=\frac{7\times136.2-21\times31.5}{196}=\frac{291.9}{196}=1.489 & 逆行列の1行めをかける \\
b=\frac{-21\times136.2+91\times31.5}{196}=\frac{6.3}{196}=0.032 & 逆行列の2行めをかける \\
\end{bkhenkei}
求めた直線は、次のとおりです。
\[ y=1.489x+0.032 \]
3-0・3-16で表計算が出した式（$a=1.4893$、$b=0.0321$）と同じです。原点を通る直線（3-16）の $a=1.497$ とは、少しちがいます。$b$ も自由に動かせるぶん、ずれの2乗の和は少し小さくなっています。
\begin{bkfigure}
\begin{tikzpicture}
\begin{axis}[bkplot, width=100mm, height=56mm, xmin=0, xmax=6.8, ymin=0, ymax=10.5, xtick={1,2,3,4,5,6}, ytick={2,4,6,8,10}, xlabel={$x$（個）}, ylabel={$y$（cm）}]
  \addplot[bkcurve, domain=0:6.6] {1.4893*x+0.0321};
  \addplot[bkGreen!60!black, line width=1pt] coordinates {(0,0) (0,0.0321)};
  \addplot[bkGreen!60!black, line width=1pt] coordinates {(1,1.6) (1,1.5214)};
  \addplot[bkGreen!60!black, line width=1pt] coordinates {(2,2.9) (2,3.0107)};
  \addplot[bkGreen!60!black, line width=1pt] coordinates {(3,4.6) (3,4.5000)};
  \addplot[bkGreen!60!black, line width=1pt] coordinates {(4,5.9) (4,5.9893)};
  \addplot[bkGreen!60!black, line width=1pt] coordinates {(5,7.6) (5,7.4786)};
  \addplot[bkGreen!60!black, line width=1pt] coordinates {(6,8.9) (6,8.9679)};
  \addplot[only marks, mark=*, mark size=1.8pt, bkOrange] coordinates {(0,0) (1,1.6) (2,2.9) (3,4.6) (4,5.9) (5,7.6) (6,8.9)};
  \node[font=\sffamily\footnotesize, text=\bkTT, anchor=north west, fill=white, inner sep=1pt] at (axis cs:3.6,4.3) {$y=1.489x+0.032$};
\end{axis}
\end{tikzpicture}
\bkfigcaption{ソラさんのデータ（オレンジ）と、正規方程式で求めた直線（青）。ずれ（残差、緑の線）は \qty{0.1}{cm} ほどしかなく、点にかくれてほとんど見えない。ずれの2乗の和が、いちばん小さくなっている。}
\end{bkfigure}

\kviifhead{表計算と Python で確かめる}
表計算では、\texttt{LINEST}（ライン・エスト）という関数で、$y=ax+b$ の $a$ と $b$ がいっぺんに出ます。A2:A8 に $x$、B2:B8 に $y$ を入れて、D2 に \texttt{=LINEST(B2:B8,A2:A8)} と入れると、D2 に $a=1.4893$、E2 に $b=0.0321$ が横にならんで出ます（はじめに $y$ の範囲、次に $x$ の範囲。確定のしかたは7-1と同じ）。Python では、正規方程式を \texttt{solve} で解いても、最小二乗法の関数 \texttt{np.linalg.lstsq} を使っても、同じ答えです。
\Needspace{11\baselineskip}
\begin{bkcode}[Python]
import numpy as np
x = np.array([0, 1, 2, 3, 4, 5, 6])
y = np.array([0, 1.6, 2.9, 4.6, 5.9, 7.6, 8.9])
A = np.column\_stack([x, np.ones(7)])    \# 1列めが x、2列めが1の行列
print(A.T @ A)                          \# A の転置×A
print(A.T @ y)
print(np.linalg.solve(A.T @ A, A.T @ y))      \# 正規方程式を解く
print(np.linalg.lstsq(A, y, rcond=None)[0])   \# 最小二乗法の関数
\end{bkcode}
\begin{bkcodeout}[実行結果]
[[91. 21.]
 [21.  7.]]
[136.2  31.5]
[1.48928571 0.03214286]
[1.48928571 0.03214286]
\end{bkcodeout}
4行めでは、$x$ と、1を7個ならべたベクトル \texttt{np.ones(7)} を、列としてならべて行列 $A$ を作りました（\texttt{column} は列、\texttt{stack} は積み重ねるという意味）。\texttt{lstsq} は least squares（最小二乗）の略です。

\begin{seikaku}
ここでは、ずれが「列のつくる面」に直角になる点が、いちばん近い点だと考えました。それは、面の上のほかの点 $\vec q$ をとると、$\vec b$ から $\vec q$ までの距離の2乗が、三平方の定理で「（ずれの長さ）$^2$＋（$A\vec x$ から $\vec q$ までの長さ）$^2$」になり、ずれの長さより必ず長くなるからです。3-16の平方完成や、のちに学ぶ偏微分（8-2）で求めても、同じ正規方程式が出てきます。また、$A$ の2本の列が平行でなければ（$x$ の値がすべて同じではなければ）、$A^{\mathsf T}A$ の行列式は0にならず、答えはただ1つに決まります。
\end{seikaku}


%% ---------------- 専門書・まとめ ----------------
\begin{senmonsho}[線形代数の教科書ふうの書き方]
\begin{bkhonbun}
$n$ 次正方行列 $A$ に対し $AX=XA=E$ を満たす $X$ が存在するとき、$A$ は正則であるといい、$X$ を $A$ の逆行列とよび $A^{-1}$ で表す。$A$ が正則であるための必要十分条件は $|A|\neq0$ である。このとき、連立1次方程式 $A\symbfit{x}=\symbfit{b}$ はただ1つの解 $\symbfit{x}=A^{-1}\symbfit{b}$ をもつ。
\end{bkhonbun}
\bkyaku{$E$ はこの本の単位行列 $I$、$|A|$ は行列式 $\det A$、太字の $\symbfit{x}$ はこの本の $\vec x$。「正則」は「逆行列がある」こと。物理第1巻（力-5）の「正方かつ正則」は「行列がたて・横同じ大きさで、逆行列がある」こと。「必要十分条件」は「同じことの言いかえ」：正則なら行列式は0でなく、行列式が0でなければ正則。}
\end{senmonsho}

\begin{taisetsu}
\def\tsl#1{\makebox[6.4\zw][l]{\sffamily\bfseries\small\color{\bkTT}#1}}
\tsl{逆行列}$\begin{pmatrix}a&b\\c&d\end{pmatrix}^{-1}=\dfrac{1}{ad-bc}\begin{pmatrix}d&-b\\-c&a\end{pmatrix}$、$A^{-1}A=AA^{-1}=I$\\[4pt]
\tsl{解き方}$A\vec x=\vec b$ の答えは $\vec x=A^{-1}\vec b$（左からかける）\\[4pt]
\tsl{順番}$(AB)^{-1}=B^{-1}A^{-1}$\\[4pt]
\tsl{掃き出し法}かける・ひく・入れかえる、の3つの操作で左を $I$ に\\[4pt]
\tsl{解けない}$\det A=0$ なら逆行列はなく、答えはないか無数\\[4pt]
\tsl{最小二乗法}正規方程式 $A^{\mathsf T}A\vec x=A^{\mathsf T}\vec b$（ずれが列に直角）
\end{taisetsu}

\begin{kotoba}
\small
\bkword{逆行列}{行列 $A$ の変換を元にもどす行列 $A^{-1}$。$A^{-1}A=AA^{-1}=I$。}
\bkword{正則}{逆行列がある（行列式が0でない）正方行列のこと。}
\bkword{掃き出し法}{係数と右辺の表に、行の3つの操作をくり返して、左を単位行列にする解き方。}
\bkword{ガウスの消去法}{ななめより下を0にしてから、下の行から代入していく解き方。コンピュータが使う。}
\bkword{階数}{掃き出したあとに残る、左が全部0ではない行の数。本当に役に立っている式の数。}
\bkword{正規方程式}{最小二乗法の答えを決める式 $A^{\mathsf T}A\vec x=A^{\mathsf T}\vec b$。}
\end{kotoba}

\bkkigocards{\bkkigocard{$A^{-1}$}{エー・インバース}{$A$ の逆行列。$A$ を元にもどす行列。$\frac{1}{A}$ とは書かない}{inverse matrix}{7-5}}{}

\begin{column}{ガウスの消去法と『九章算術』}
コンピュータが使う「ガウスの消去法」には、ドイツの数学者ガウスの名前がついています。ガウスは、星の動きを調べる研究の中で、最小二乗法の正規方程式を、この消去法で解く手順をまとめました（1809年の本などで）。でも、係数を表に並べて、行どうしをひいて文字を消していく方法そのものは、7-1のコラムで見た中国の『九章算術』（紀元1世紀ごろまでにまとめられた）の「方程」の章に、すでに書かれていました。2000年前の人たちが算木で計算していた手順が、いまでは、橋や飛行機の設計で、何百万個もの式をコンピュータで解く手順になっているのです。
\end{column}

%% ---------------- 例題 ----------------
\bksection{reidai}{例題でたしかめよう}

\begin{reidai}[2×2 の逆行列を作って確かめる]
次の行列の逆行列を作り、もとの行列にかけて $I$ になることを確かめましょう。
\[ \text{（1）}\ \begin{pmatrix}3&1\\5&2\end{pmatrix}\qquad\qquad \text{（2）}\ \begin{pmatrix}4&2\\2&4\end{pmatrix} \]
\kangaekata ①の3つの手順：$a$ と $d$ を入れかえ、$b$ と $c$ にマイナス、行列式でわる。\par
\shiki（1）行列式 $3\times2-1\times5=1$。逆行列 $\begin{pmatrix}2&-1\\-5&3\end{pmatrix}$\kviifq
確かめ：$\begin{pmatrix}2&-1\\-5&3\end{pmatrix}\begin{pmatrix}3&1\\5&2\end{pmatrix}=\begin{pmatrix}6-5&2-2\\-15+15&-5+6\end{pmatrix}=I$\kviifq
（2）行列式 $4\times4-2\times2=12$。逆行列 $\dfrac{1}{12}\begin{pmatrix}4&-2\\-2&4\end{pmatrix}$\kviifq
確かめ：$\dfrac{1}{12}\begin{pmatrix}16-4&8-8\\-8+8&-4+16\end{pmatrix}=\dfrac{1}{12}\begin{pmatrix}12&0\\0&12\end{pmatrix}=I$\par
\kotae（1）\nobreak $\begin{pmatrix}2&-1\\-5&3\end{pmatrix}$　（2）\nobreak $\dfrac{1}{12}\begin{pmatrix}4&-2\\-2&4\end{pmatrix}$
\end{reidai}

\begin{jibun}
次の行列の逆行列を作り、もとの行列にかけて $I$ になることを確かめましょう。
\[ \text{（1）}\ \begin{pmatrix}5&2\\2&1\end{pmatrix}\qquad\qquad \text{（2）}\ \begin{pmatrix}2&1\\1&3\end{pmatrix} \]
\bkjans{（1）\nobreak 行列式1。逆行列は1行め $(1,\ -2)$、2行め $(-2,\ 5)$。（2）\nobreak 行列式5。逆行列は $\frac{1}{5}$ に1行め $(3,\ -1)$、2行め $(-1,\ 2)$をかけたもの。どちらも、もとの行列にかけると $I$。}
\end{jibun}

\begin{reidai}[逆行列で連立方程式を解く]
次の連立方程式を、逆行列を使って解きましょう。
\[ \begin{cases}2x+3y=420\\2x+y=300\end{cases} \]
\kangaekata $A\vec x=\vec b$ と書いて、$\vec x=A^{-1}\vec b$。逆行列は左からかける。\par
\shiki $A=\begin{pmatrix}2&3\\2&1\end{pmatrix}$、行列式 $2\times1-3\times2=-4$\kviifq
$A^{-1}=\dfrac{1}{-4}\begin{pmatrix}1&-3\\-2&2\end{pmatrix}$\kviifq
$\vec x=\dfrac{1}{-4}\begin{pmatrix}1\times420-3\times300\\-2\times420+2\times300\end{pmatrix}=\dfrac{1}{-4}\begin{pmatrix}-480\\-240\end{pmatrix}=\begin{pmatrix}120\\60\end{pmatrix}$\kviifq
確かめ：$2\times120+3\times60=420$、$2\times120+60=300$\par
\kotae $x=120$、$y=60$
\end{reidai}

\begin{jibun}
次の連立方程式を、逆行列を使って解きましょう。
\[ \begin{cases}3x+2y=12\\x+y=5\end{cases} \]
\bkjans{行列式 $3-2=1$。$A^{-1}$ は1行め $(1,\ -2)$、2行め $(-1,\ 3)$。$x=12-10=2$、$y=-12+15=3$。}
\end{jibun}

\begin{reidai}[同じ逆行列で、右辺だけを変える]
次の回路の式（電圧 $E_1$、$E_2$〔V〕、電流 $I_1$、$I_2$〔A〕）で、（1）$E_1=9$、$E_2=6$、（2）$E_1=12$、$E_2=0$ のときの電流を求めましょう。
\[ \begin{cases}4I_1+2I_2=E_1\\2I_1+4I_2=E_2\end{cases} \]
\kangaekata 逆行列（$\dfrac{1}{12}$ に1行め $(4,\ -2)$、2行め $(-2,\ 4)$ をかけたもの）は電圧に関係しないので、1回作って両方に使う。\par
\shiki（1）$\dfrac{1}{12}\begin{pmatrix}4\times9-2\times6\\-2\times9+4\times6\end{pmatrix}=\dfrac{1}{12}\begin{pmatrix}24\\6\end{pmatrix}=\begin{pmatrix}2\\0.5\end{pmatrix}$\kviifq
（2）$\dfrac{1}{12}\begin{pmatrix}4\times12-2\times0\\-2\times12+4\times0\end{pmatrix}=\dfrac{1}{12}\begin{pmatrix}48\\-24\end{pmatrix}=\begin{pmatrix}4\\-2\end{pmatrix}$\par
\kotae（1）\nobreak $I_1=\qty{2}{A}$、$I_2=\qty{0.5}{A}$。（2）\nobreak $I_1=\qty{4}{A}$、$I_2=\qty{-2}{A}$（$I_2$ は矢印と反対向きに \qty{2}{A}）
\end{reidai}

\begin{jibun}
次の回路の式で、（1）$E_1=8$、$E_2=7$、（2）$E_1=5$、$E_2=0$ のときの電流を求めましょう。
\[ \begin{cases}3I_1+2I_2=E_1\\2I_1+3I_2=E_2\end{cases} \]
\bkjans{逆行列は $\frac{1}{5}$ に1行め $(3,\ -2)$、2行め $(-2,\ 3)$をかけたもの。（1）\nobreak $\frac{1}{5}(24-14,\ -16+21)=(2,\ 1)$。$I_1=\qty{2}{A}$、$I_2=\qty{1}{A}$。（2）\nobreak $\frac{1}{5}(15,\ -10)=(3,\ -2)$。$I_1=\qty{3}{A}$、$I_2=\qty{-2}{A}$。}
\end{jibun}

\begin{reidai}[つないだ機械の逆行列]
次の $A$、$B$ について、$(AB)^{-1}$ を2通りで求めて、くらべましょう。
\[ A=\begin{pmatrix}2&1\\1&1\end{pmatrix},\qquad B=\begin{pmatrix}1&0\\2&1\end{pmatrix} \]\par
\kangaekata（ア）$AB$ を計算してから逆行列。（イ）$B^{-1}A^{-1}$（順番が逆）。\par
\shiki（ア）$AB=\begin{pmatrix}2+2&0+1\\1+2&0+1\end{pmatrix}=\begin{pmatrix}4&1\\3&1\end{pmatrix}$、行列式 $4-3=1$\kviifq
$(AB)^{-1}=\begin{pmatrix}1&-1\\-3&4\end{pmatrix}$\kviifq
（イ）$A^{-1}=\begin{pmatrix}1&-1\\-1&2\end{pmatrix}$、$B^{-1}=\begin{pmatrix}1&0\\-2&1\end{pmatrix}$（どちらも行列式1）\kviifq
$B^{-1}A^{-1}=\begin{pmatrix}1&-1\\-2-1&2+2\end{pmatrix}=\begin{pmatrix}1&-1\\-3&4\end{pmatrix}$\par
\kotae どちらも $\begin{pmatrix}1&-1\\-3&4\end{pmatrix}$
\end{reidai}

\begin{jibun}
次の $A$、$B$ について、$(AB)^{-1}$ を2通りで求めて、くらべましょう。
\[ A=\begin{pmatrix}3&2\\1&1\end{pmatrix},\qquad B=\begin{pmatrix}1&1\\0&1\end{pmatrix} \]
\bkjans{（ア）\nobreak $AB$ は1行め $(3,\ 5)$、2行め $(1,\ 2)$、行列式1、$(AB)^{-1}$ は1行め $(2,\ -5)$、2行め $(-1,\ 3)$。（イ）\nobreak $B^{-1}$ は1行め $(1,\ -1)$、2行め $(0,\ 1)$、$A^{-1}$ は1行め $(1,\ -2)$、2行め $(-1,\ 3)$、$B^{-1}A^{-1}$ は1行め $(2,\ -5)$、2行め $(-1,\ 3)$。同じ。}
\end{jibun}

\begin{reidai}[掃き出し法で3つの文字]
次の連立方程式を、掃き出し法で解きましょう。
\[ \begin{cases}y+3z=11\\x+y+z=6\\x-y+2z=5\end{cases} \]
\kangaekata 1行め1列めが0なので、このままでは1にできない。1列めが0でない②と入れかえてから（操作3）、1列めから順に、ななめを1、ほかを0にする。右辺もいっしょに計算する。\par
\shiki $\kviifaug{rrr|r}{0&1&3&11\\1&1&1&6\\1&-1&2&5}\xrightarrow{\text{①}\leftrightarrow\text{②}}\kviifaug{rrr|r}{1&1&1&6\\0&1&3&11\\1&-1&2&5}$\kviifq
$\xrightarrow{\text{③}-\text{①}}\kviifaug{rrr|r}{1&1&1&6\\0&1&3&11\\0&-2&1&-1}\xrightarrow[\text{③}+2\times\text{②}]{\text{①}-\text{②}}\kviifaug{rrr|r}{1&0&-2&-5\\0&1&3&11\\0&0&7&21}$\kviifq
$\xrightarrow{\text{③}\div7}\kviifaug{rrr|r}{1&0&-2&-5\\0&1&3&11\\0&0&1&3}\xrightarrow[\text{②}-3\times\text{③}]{\text{①}+2\times\text{③}}\kviifaug{rrr|r}{1&0&0&1\\0&1&0&2\\0&0&1&3}$\kviifq
確かめ：$2+9=11$、$1+2+3=6$、$1-2+6=5$\par
\kotae $x=1$、$y=2$、$z=3$
\end{reidai}

\begin{jibun}
次の連立方程式を、掃き出し法で解きましょう。
\[ \begin{cases}x+2y+z=8\\x+y-z=0\\2x-y+z=3\end{cases} \]
\bkjans{②－①、③－2×① で2行め $(0,\ -1,\ -2\mathrel{|}-8)$、3行め $(0,\ -5,\ -1\mathrel{|}-13)$。②×$(-1)$ で $(0,\ 1,\ 2\mathrel{|}8)$。①－2×②、③＋5×② で1行め $(1,\ 0,\ -3\mathrel{|}-8)$、3行め $(0,\ 0,\ 9\mathrel{|}27)$。③÷9 で $z=3$。①＋3×③、②－2×③ で $x=1$、$y=2$。}
\end{jibun}

\begin{reidai}[掃き出し法で逆行列]
次の行列の逆行列を、掃き出し法で求めましょう。また、①の3つの手順の答えとくらべましょう。
\[ \begin{pmatrix}1&2\\3&7\end{pmatrix} \]\par
\kangaekata 右に $I$ をならべて、左が $I$ になるまで掃き出す。右に出てきた行列が逆行列。\par
\shiki $\kviifaug{rr|rr}{1&2&1&0\\3&7&0&1}\xrightarrow{\text{②}-3\times\text{①}}\kviifaug{rr|rr}{1&2&1&0\\0&1&-3&1}$\kviifq
$\xrightarrow{\text{①}-2\times\text{②}}\kviifaug{rr|rr}{1&0&7&-2\\0&1&-3&1}$\kviifq
3つの手順：行列式 $7-6=1$。$\begin{pmatrix}7&-2\\-3&1\end{pmatrix}$\par
\kotae $\begin{pmatrix}7&-2\\-3&1\end{pmatrix}$（3つの手順の答えと同じ）
\end{reidai}

\begin{jibun}
次の行列の逆行列を、掃き出し法で求めましょう。また、①の3つの手順の答えとくらべましょう。
\[ \begin{pmatrix}1&3\\2&7\end{pmatrix} \]
\bkjans{②－2×① で $(0,\ 1\mathrel{|}-2,\ 1)$。①－3×② で $(1,\ 0\mathrel{|}7,\ -3)$。逆行列は1行め $(7,\ -3)$、2行め $(-2,\ 1)$。行列式 $7-6=1$ の3つの手順と同じ。}
\end{jibun}

\begin{reidai}[答えがあるかどうかを調べる]
次の連立方程式を掃き出し法で調べ、答えがないか、無数にあるかを答えましょう。\par
\[ \text{（1）}\ \begin{cases}x+2y=3\\2x+4y=6\end{cases}\qquad\qquad \text{（2）}\ \begin{cases}x+2y=3\\2x+4y=5\end{cases} \]
\kangaekata 行列式は $4-4=0$。②－2×① で2行めの左が0になる。右が0なら無数、0でなければ答えはない。\par
\shiki（1）②－2×①：$(2-2,\ 4-4\mathrel{|}6-6)=(0,\ 0\mathrel{|}0)$。「$0=0$」\kviifq
（2）②－2×①：$(2-2,\ 4-4\mathrel{|}5-6)=(0,\ 0\mathrel{|}-1)$。「$0=-1$」\par
\kotae（1）\nobreak 答えは無数（$x=3-2y$ なら何でもよい）。（2）\nobreak 答えはない
\end{reidai}

\begin{jibun}
次の連立方程式を掃き出し法で調べ、答えがないか、無数にあるかを答えましょう。\par
\[ \text{（1）}\ \begin{cases}3x-y=2\\6x-2y=4\end{cases}\qquad\qquad \text{（2）}\ \begin{cases}3x-y=2\\6x-2y=1\end{cases} \]
\bkjans{（1）\nobreak ②－2×① で $(0,\ 0\mathrel{|}0)$。答えは無数（$y=3x-2$）。（2）\nobreak ②－2×① で $(0,\ 0\mathrel{|}-3)$、「$0=-3$」。答えはない。}
\end{jibun}

\begin{reidai}[ロボットの関節の角速度]
2つの関節のロボットの腕で、ヤコビ行列が次のとき（$\det J=0.0530$）、
\[ J=\begin{pmatrix}-0.3915&-0.2415\\0.3245&0.0647\end{pmatrix} \]
手をまっすぐ上へ \qty{0.1}{m/s} で動かす（手の速度 $(0,\ 0.1)$）には、肩とひじを何 rad/s で回せばよいですか。\par
\kangaekata 手の速度 $(0,\ 0.1)$ を右辺にして、逆行列で解く。右辺の1行めが0なので、逆行列の2列めだけが効く。\par
\shiki 逆行列：$J^{-1}=\dfrac{1}{0.0530}\begin{pmatrix}0.0647&0.2415\\-0.3245&-0.3915\end{pmatrix}$\kviifq
$\dot{\vec\theta}=\dfrac{1}{0.0530}\begin{pmatrix}0.2415\times0.1\\-0.3915\times0.1\end{pmatrix}=\dfrac{1}{0.0530}\begin{pmatrix}0.02415\\-0.03915\end{pmatrix}$\kviifq
$\dfrac{1}{0.0530}\begin{pmatrix}0.02415\\-0.03915\end{pmatrix}=\begin{pmatrix}0.456\\-0.739\end{pmatrix}$\par
\kotae 肩を約 \qty{0.46}{rad/s}（左回り）、ひじを約 \qty{0.74}{rad/s}（右回り）
\end{reidai}
\noindent 丸める前の値で計算すると0.455と $-0.738$ です（小さなちがいは丸めのため）。確かめに $J$ をかけると、ほしい速度になります。
\[ -0.3915\times0.456+0.2415\times0.739=0.000 \]
\[ 0.3245\times0.456-0.0647\times0.739=0.100 \]

\begin{jibun}
例題8と同じロボットの腕で、手をまっすぐ左へ \qty{0.1}{m/s} で動かす（手の速度 $(-0.1,\ 0)$）には、肩とひじを何 rad/s で回せばよいですか。
\bkjans{肩：$0.0647\times(-0.1)\div0.0530\approx-0.122$。ひじ：$-0.3245\times(-0.1)\div0.0530\approx0.612$。肩を約 \qty{0.12}{rad/s}（右回り）、ひじを約 \qty{0.61}{rad/s}（左回り）。}
\end{jibun}

\begin{reidai}[最小二乗法で直線を引く]
4つの点 $(1,\ 3)$、$(2,\ 5)$、$(3,\ 6)$、$(4,\ 8)$ に、最小二乗法で直線 $y=ax+b$ を引きましょう。\par
\kangaekata 正規方程式（左の行列は「$x^2$ の和、$x$ の和／$x$ の和、点の数」、右辺は「$xy$ の和／$y$ の和」）を作って、逆行列で解く。\par
\shiki $x$ の和 $1+2+3+4=10$、$x^2$ の和 $1+4+9+16=30$\kviifq
$y$ の和 $3+5+6+8=22$、$xy$ の和 $3+10+18+32=63$\kviifq
正規方程式 $\begin{pmatrix}30&10\\10&4\end{pmatrix}\begin{pmatrix}a\\b\end{pmatrix}=\begin{pmatrix}63\\22\end{pmatrix}$、行列式 $120-100=20$\kviifq
$a=\dfrac{4\times63-10\times22}{20}=\dfrac{32}{20}=1.6$、$b=\dfrac{-10\times63+30\times22}{20}=\dfrac{30}{20}=1.5$\par
\kotae $y=1.6x+1.5$
\end{reidai}

\begin{jibun}
おもり $x$ 個が1〜5のとき、ばねの長さ $y$〔cm〕は14.0、15.9、18.1、20.1、21.9でした（3-16のばねC）。最小二乗法で、直線 $y=ax+b$ を求めましょう。
\bkjans{$x$ の和15、$x^2$ の和55、$y$ の和90.0、$xy$ の和290.0。正規方程式の行列は1行め $(55,\ 15)$、2行め $(15,\ 5)$、右辺は $(290,\ 90)$、行列式 $275-225=50$。$a=(5\times290-15\times90)\div50=2$、$b=(-15\times290+55\times90)\div50=12$。$y=2x+12$（3-16と同じ）。}
\end{jibun}

%% ---------------- まちがえやすいところ ----------------
\bksection{tsumazuki}{まちがえやすいところをチェック}

\begin{tsumazuki}{逆行列をかける場所}
\bkng{$A\vec x=\vec b$ だから $\vec x=\vec bA^{-1}$}
\bkok{$A\vec x=\vec b$ だから $\vec x=A^{-1}\vec b$（左からかける）}
$A$ は $\vec x$ の左にいるので、消すための $A^{-1}$ も左に置きます。$\vec bA^{-1}$ は、2行1列×2行2列で、大きさが合わず計算できません。
\end{tsumazuki}

\begin{tsumazuki}{逆行列の作り方}
\bkng{$\begin{pmatrix}2&3\\2&1\end{pmatrix}^{-1}=\begin{pmatrix}2&-3\\-2&1\end{pmatrix}$}
\bkok{$\begin{pmatrix}2&3\\2&1\end{pmatrix}^{-1}=\dfrac{1}{-4}\begin{pmatrix}1&-3\\-2&2\end{pmatrix}$}
左上と右下を入れかえるのを忘れ、行列式でわるのも忘れています。作ったら、もとの行列にかけて $I$ になるか確かめましょう。まちがいにすぐ気づけます。
\end{tsumazuki}

\begin{tsumazuki}{つないだ機械の逆行列}
\bkng{$(AB)^{-1}=A^{-1}B^{-1}$}
\bkok{$(AB)^{-1}=B^{-1}A^{-1}$（順番が逆）}
あとからしたことを、先にもどします（くつ下とくつ）。行列ではかける順で答えが変わるので、順番をまちがえると、ちがう行列になります。
\end{tsumazuki}

\begin{tsumazuki}{行列式が0のとき}
\bkng{$\begin{pmatrix}1&2\\2&4\end{pmatrix}$ の逆行列を、3つの手順で作る}
\bkok{行列式が $4-4=0$ なので、逆行列はない}
0でわることになるので、逆行列は作れません。表計算は \texttt{\#NUM!}、Python は \texttt{LinAlgError} と出ます。はじめに行列式を計算するくせをつけましょう。
\end{tsumazuki}

\begin{tsumazuki}{右辺をひき忘れる}
\bkng{②－① で、左の係数だけをひいて、右辺はそのまま}
\bkok{右辺もいっしょにひく（1行全体をひく）}
1つの行は1つの式です。式から式をひくときは、左辺も右辺もひきます。右辺を忘れると、ちがう式になってしまいます。
\end{tsumazuki}
\begin{tsumazuki}{0の所で、入れかえない}
\bkng{1行め1列めが0なのに、①÷0 をしようとする（または、行ではなく列を入れかえる）}
\bkok{1列めが0でない行と、行ごと入れかえる（右辺もいっしょに）}
0ではわれません。行を入れかえるのは、式を書く順番を変えるだけなので、答えは変わりません。列を入れかえると、$x$ と $y$ の場所が入れかわるので、答えを読むときにまちがえます。
\end{tsumazuki}

\begin{tsumazuki}{正規方程式のかける順}
\bkng{$AA^{\mathsf T}\vec x=A^{\mathsf T}\vec b$}
\bkok{$A^{\mathsf T}A\vec x=A^{\mathsf T}\vec b$}
ソラさんのデータでは、$A$ は7行2列です。$A^{\mathsf T}A$ は2行2列で、$a$ と $b$ の2つの式になります。$AA^{\mathsf T}$ は7行7列になり、$\vec x$（2行1列）にかけられません。「$A\vec x=\vec b$ に、左から $A^{\mathsf T}$ をかけた形」と覚えましょう。
\end{tsumazuki}
\begin{tsumazuki}{逆行列を $\frac{1}{A}$ と書く}
\bkng{$\vec x=\dfrac{\vec b}{A}$}
\bkok{$\vec x=A^{-1}\vec b$}
行列のわり算は決めていません。$\dfrac{\vec b}{A}$ では、$A^{-1}$ を左からかけるのか右からかけるのかがわかりません。
\end{tsumazuki}

%% ---------------- 物理・工学 ----------------
\begin{butsuri}
\begin{itemize}[leftmargin=1.2\zw, itemsep=3pt, topsep=0pt, label={\color{bkPurple}●}]
\item \textbf{橋やビルの設計：何百万個の連立方程式}\par\noindent\begin{minipage}[c]{0.70\linewidth}
橋やビルを設計するときは、建物を小さな部品（要素）に細かく分け、部品どうしがつながる点ごとに、力のつり合いの式を立てます。これを「有限要素法」といいます。7-5の3つのばねの $K\vec u=\vec F$ と同じ形の連立方程式で、大きな建物では式が何百万個にもなります（物理 力-16）。それを、この章のガウスの消去法の考えを使ったプログラムで解いて、どこにどれだけの力がかかるかを調べています。
\end{minipage}\hfill
\begin{minipage}[c]{0.27\linewidth}\centering
\begin{tikzpicture}[x=1mm, y=1mm]
  \foreach \kviifi in {0,...,5} {\draw[bkInk!60, line width=0.4pt] (\kviifi*5,0) -- (\kviifi*5,15);}
  \foreach \kviifj in {0,...,3} {\draw[bkInk!60, line width=0.4pt] (0,\kviifj*5) -- (25,\kviifj*5);}
  \foreach \kviifi in {0,...,4} {\foreach \kviifj in {0,...,2} {\draw[bkInk!40, line width=0.3pt] (\kviifi*5,\kviifj*5) -- (\kviifi*5+5,\kviifj*5+5);}}
  \foreach \kviifi in {0,...,5} {\foreach \kviifj in {0,...,3} {\fill[bkRed] (\kviifi*5,\kviifj*5) circle (0.5);}}
\end{tikzpicture}
\end{minipage}
\item \textbf{電子回路のシミュレーション}\par\noindent 電子回路の設計では、回路の各点の電圧をわからない数として、キルヒホッフの法則（3-6のきまり1〜3）の式を立て、コンピュータで解きます（物理 電-5）。トランジスタのように、電流と電圧が比例しない部品があるときも、少しずつ直しながら、何度も連立1次方程式を解いて答えに近づけます。
\item \textbf{ロボットの腕：ヤコビ行列の逆行列}\par\noindent 工場のロボットは、手を決まった道すじで動かすために、手の速度 $\dot{\vec p}$ から関節の角速度を $\dot{\vec\theta}=J^{-1}\dot{\vec p}$ で、1秒に何百回も計算しています（例題8）。腕がのびきる近くでは $\det J$ が0に近づき、関節を急に速く回さなければならなくなります（章末問題14）。そこで、その近くを通らないように、道すじを決めておくのです。
\end{itemize}
\end{butsuri}

\begin{butsuri}
\begin{itemize}[leftmargin=1.2\zw, itemsep=3pt, topsep=0pt, label={\color{bkPurple}●}]
\item \textbf{GPS：衛星までの距離から、位置を決める}\par\noindent\begin{minipage}[c]{0.70\linewidth}
スマホの GPS は、空の人工衛星から届く電波の時刻から、衛星までの距離を求めます。わからない数は、自分の位置 $x$、$y$、$z$ と、スマホの時計のずれの4つなので、4つ以上の衛星が必要です。距離の式は2乗をふくむので、近い位置を仮に決め、そこからの「少しのずれ」についての連立1次方程式を解いて、位置を直すことをくり返します。5つ以上の衛星が見えるときは、式が多すぎるので、最小二乗法（⑦）でいちばんよく合う位置を求めます。
\end{minipage}\hfill
\begin{minipage}[c]{0.27\linewidth}\centering
\begin{tikzpicture}[x=1mm, y=1mm]
  \draw[bkGreen!60!black, line width=1pt] (0,0) arc (180:0:14 and 4);
  \fill[bkInk] (14,1.5) circle (1);
  \foreach \kviifx/\kviify in {2/22, 12/26, 22/23, 27/14} {\filldraw[fill=bkGray!40, draw=bkInk!70, line width=0.5pt] (\kviifx-1.5,\kviify-1) rectangle (\kviifx+1.5,\kviify+1); \draw[bkInk!50, line width=0.4pt, dashed] (\kviifx,\kviify-1) -- (14,2.5);}
\end{tikzpicture}
\end{minipage}
\item \textbf{次元解析：単位の指数を逆行列で}\par\noindent 物理 力-15では、長さ・質量・時間の単位の指数を並べた行列の行列式が1（7-4の章末問題）なので、逆行列があり、ほしい量を作る指数の組がただ1つに決まることを使いました。行列式が0でない（正則）ことが、答えがただ1つに決まる保証になっているのです。
\item \textbf{センサーの目もり合わせ（校正）}\par\noindent 温度センサーや圧力センサーは、はかった量を電圧に変えて出します。工場では、温度がわかっている所でセンサーの電圧を何回もはかり、「温度と電圧の直線」を最小二乗法（⑦）で決めておきます。こうして決めた直線を使って、電圧から温度を計算するのです。はかる点の数が何十個でも、解くのは2つの文字の正規方程式だけです。
\end{itemize}
\end{butsuri}

%% ---------------- 章末問題 ----------------
\bkendproblems
\bkproblemgroup{基本}{4}

\bkproblem 次の行列の逆行列を作り、もとの行列にかけて $I$ になることを確かめましょう。
\[ \text{（1）}\ \begin{pmatrix}4&3\\1&1\end{pmatrix}\qquad\qquad \text{（2）}\ \begin{pmatrix}3&1\\1&3\end{pmatrix} \]
\bkpans{（1）\nobreak 行列式1。逆行列は1行め $(1,\ -3)$、2行め $(-1,\ 4)$。（2）\nobreak 行列式8。逆行列は $\frac{1}{8}$ に1行め $(3,\ -1)$、2行め $(-1,\ 3)$をかけたもの。}

\bkproblem 次の連立方程式を、逆行列を使って解きましょう。
\[ \begin{cases}5x+2y=16\\2x+y=7\end{cases} \]
\bkpans{行列式 $5-4=1$。$A^{-1}$ は1行め $(1,\ -2)$、2行め $(-2,\ 5)$。$x=16-14=2$、$y=-32+35=3$。}

\bkproblem 次の回路の式で、（1）$E_1=12$、$E_2=9$、（2）$E_1=21$、$E_2=0$ のときの電流 $I_1$、$I_2$〔A〕を求めましょう。
\[ \begin{cases}5I_1+2I_2=E_1\\2I_1+5I_2=E_2\end{cases} \]
\bkpans{逆行列は $\frac{1}{21}$ に1行め $(5,\ -2)$、2行め $(-2,\ 5)$をかけたもの。（1）\nobreak $\frac{1}{21}(60-18,\ -24+45)=(2,\ 1)$。（2）\nobreak $\frac{1}{21}(105,\ -42)=(5,\ -2)$。}

\bkproblem 次の $A$、$B$ について、$(AB)^{-1}$ を2通り（$AB$ の逆行列と、$B^{-1}A^{-1}$）で求めましょう。
\[ A=\begin{pmatrix}1&1\\1&2\end{pmatrix},\qquad B=\begin{pmatrix}2&1\\1&1\end{pmatrix} \]
\bkpans{$AB$ は1行め $(3,\ 2)$、2行め $(4,\ 3)$、行列式1、$(AB)^{-1}$ は1行め $(3,\ -2)$、2行め $(-4,\ 3)$。$B^{-1}$ は1行め $(1,\ -1)$、2行め $(-1,\ 2)$、$A^{-1}$ は1行め $(2,\ -1)$、2行め $(-1,\ 1)$、$B^{-1}A^{-1}$ は1行め $(3,\ -2)$、2行め $(-4,\ 3)$。}

\bkproblem 次の連立方程式を、掃き出し法で解きましょう。
\[ \begin{cases}x+y+z=6\\x+2y-z=2\\2x-y+z=3\end{cases} \]
\bkpans{②－①、③－2×① で $(0,\ 1,\ -2\mathrel{|}-4)$、$(0,\ -3,\ -1\mathrel{|}-9)$。①－②、③＋3×② で $(1,\ 0,\ 3\mathrel{|}10)$、$(0,\ 0,\ -7\mathrel{|}-21)$。③÷$(-7)$ で $z=3$。$x=10-9=1$、$y=-4+6=2$。}

\bkproblem 次の連立方程式を、掃き出し法で解きましょう（例題5）。
\[ \begin{cases}y+2z=8\\x+y+z=6\\x+2y=5\end{cases} \]
\bkpans{1行め1列めが0なので、①と②を入れかえて1行め $(1,\ 1,\ 1\mathrel{|}6)$、2行め $(0,\ 1,\ 2\mathrel{|}8)$。③－① で $(0,\ 1,\ -1\mathrel{|}-1)$。①－②、③－② で $(1,\ 0,\ -1\mathrel{|}-2)$、$(0,\ 0,\ -3\mathrel{|}-9)$。③÷$(-3)$ で $z=3$。$x=-2+3=1$、$y=8-6=2$。}

\bkproblem 次の行列の逆行列を、掃き出し法で求めましょう。
\[ \begin{pmatrix}1&2\\2&5\end{pmatrix} \]
\bkpans{②－2×① で $(0,\ 1\mathrel{|}-2,\ 1)$。①－2×② で $(1,\ 0\mathrel{|}5,\ -2)$。逆行列は1行め $(5,\ -2)$、2行め $(-2,\ 1)$。}

\bkproblem 次の連立方程式を掃き出し法で調べ、答えがないか、無数にあるかを答えましょう。\par
\[ \text{（1）}\ \begin{cases}2x+6y=4\\x+3y=2\end{cases}\qquad\qquad \text{（2）}\ \begin{cases}2x+6y=4\\x+3y=1\end{cases} \]
\bkpans{②－0.5×① で、2行めの左を0にする（例題7と同じ考え方）。（1）\nobreak $(1-1,\ 3-3\mathrel{|}2-2)=(0,\ 0\mathrel{|}0)$ で答えは無数（$x=2-3y$）。（2）\nobreak $(1-1,\ 3-3\mathrel{|}1-2)=(0,\ 0\mathrel{|}-1)$、「$0=-1$」で答えはない。}

\bkproblemgroup{チャレンジ}{4}

\bkproblem 4つの点 $(0,\ 1)$、$(1,\ 2)$、$(2,\ 4)$、$(3,\ 5)$ に、最小二乗法で直線 $y=ax+b$ を引きましょう。
\bkpans{$x$ の和6、$x^2$ の和14、$y$ の和12、$xy$ の和25。正規方程式の行列は1行め $(14,\ 6)$、2行め $(6,\ 4)$、右辺は $(25,\ 12)$、行列式 $56-36=20$。$a=(4\times25-6\times12)\div20=1.4$、$b=(-6\times25+14\times12)\div20=0.9$。$y=1.4x+0.9$。}

\bkproblem 例題8のロボットの腕（$\det J=0.0530$）で、手の速度を $(0.1,\ 0.1)$（m/s）にするには、肩とひじを何 rad/s で回せばよいですか（②の本文の $(0.1,\ 0)$ の答えと、例題8の答えをたしてもよい）。
\bkpans{肩：$0.0647+0.2415=0.3062$、$0.3062\times0.1\div0.0530\approx0.578$。ひじ：$-0.3245-0.3915=-0.716$、$-0.716\times0.1\div0.0530\approx-1.351$。肩約 \qty{0.58}{rad/s}、ひじ約 \qty{-1.35}{rad/s}（$(0.122,\ -0.612)$ と $(0.456,\ -0.739)$ の和とほぼ同じ）。}

\bkproblem 4-11の看板で、左のひもが水平から45°、右のひもも45°のとき、重さ \qty{100}{N} をつるすひもの力 $T_1$、$T_2$ を、逆行列で求めましょう（$\cos45°$ も $\sin45°$ も0.707とする）。
\bkpans{行列は1行め $(-0.707,\ 0.707)$、2行め $(0.707,\ 0.707)$。行列式は $-0.5-0.5=-1$。逆行列は、もとの行列と同じ。$T_1=\qty{70.7}{N}$、$T_2=\qty{70.7}{N}$。}

\bkproblem 次の連立方程式を、掃き出し法で解きましょう。
\[ \begin{cases}x+y=3\\y+z=5\\x+z=4\end{cases} \]
\bkpans{表は1行め $(1,\ 1,\ 0\mathrel{|}3)$、2行め $(0,\ 1,\ 1\mathrel{|}5)$、3行め $(1,\ 0,\ 1\mathrel{|}4)$。③－① で $(0,\ -1,\ 1\mathrel{|}1)$。③＋② で $(0,\ 0,\ 2\mathrel{|}6)$、$z=3$。②から $y=5-3=2$、①から $x=3-2=1$。}

\bkproblem 3つの点 $(0,\ 2)$、$(1,\ 2)$、$(2,\ 5)$ に、最小二乗法で直線 $y=ax+b$ を引きましょう。また、ずれの和が0になることを確かめましょう。
\bkpans{$x$ の和3、$x^2$ の和5、$y$ の和9、$xy$ の和12。正規方程式の行列は1行め $(5,\ 3)$、2行め $(3,\ 3)$、右辺は $(12,\ 9)$、行列式6。$a=(3\times12-3\times9)\div6=1.5$、$b=(-3\times12+5\times9)\div6=1.5$。$y=1.5x+1.5$。ずれは $0.5$、$-1$、$0.5$ で、和は0。}

\bkproblem ひじの角が10°で、腕がのびきる近くのロボットの腕では、ヤコビ行列は次のとおりで、$\det J=0.0130$ です。手をまっすぐ上へ \qty{0.1}{m/s} で動かす（手の速度 $(0,\ 0.1)$）には、肩とひじを何 rad/s で回せばよいですか。例題8（ひじ45°）の答えとくらべましょう。
\[ J=\begin{pmatrix}-0.3107&-0.1607\\0.4513&0.1915\end{pmatrix} \]
\bkpans{肩：$0.1607\times0.1\div0.0130\approx1.23$。ひじ：$-0.3107\times0.1\div0.0130\approx-2.39$。肩約 \qty{1.23}{rad/s}、ひじ約 \qty{-2.39}{rad/s}。例題8（0.46、$-0.74$）の約2.7倍・約3.2倍の速さがいる。}

\begin{furikaeri}
\begin{checklist}
\item 2×2 の逆行列を作って確かめ、$\vec x=A^{-1}\vec b$ で連立方程式を解ける。\bkfurref{bk:7-5:gyaku}
\item 掃き出し法で、3つの文字の連立方程式と逆行列を求められる。\bkfurref{bk:7-5:hakidashi}
\item 正規方程式 $A^{\mathsf T}A\vec x=A^{\mathsf T}\vec b$ で、最小二乗法の直線を求められる。\bkfurref{bk:7-5:saisho}
\end{checklist}
\end{furikaeri}

\bkfukidashi{sora}{逆行列と掃き出し法で、式が何個あっても解けるようになった。第2部では、行列で「回す」ことを考えるよ。}
